\documentclass[11pt]{article}
\usepackage[a4paper,margin=17mm]{geometry}
\usepackage[no-math]{fontspec}
\setmainfont{Latin Modern Roman}
\usepackage{amsmath,amssymb,amsthm,xfrac,xspace,hyperref,graphicx,comment}
\excludecomment{DefTralics}
\providecommand{\dpl}{\displaystyle}
\providecommand{\backmatter}{}
\providecommand{\loccit}{\emph{loc. cit.}}

\providecommand{\bibliofont}{\small}
\newcommand{\ie}{i.e.,\xspace}
\newcommand{\eg}{e.g.,\xspace}
\newcommand{\etc}{etc.\xspace}
\newcommand{\cf}{cf.\xspace}
\newcommand{\resp}{resp.\xspace}
\newcommand{\smaller}{\small}
\newcommand{\Subsection}{\subsection}
\newcommand{\Subsubsection}{\subsubsection}
\newcommand{\skpt}{\smallskip}
\emergencystretch=3em
\providecommand{\skippointrait}{\smallskip}
\renewcommand{\proofname}{Démonstration}

% Begin literal retained input author-preamble.tex
\usepackage[scr=rsfs,cal=euler]{mathalfa}
\multlinegap0pt
\hyphenation{admis-sible affine con-nexe con-nexes dimen-sion domai-ne domai-nes géné-ral obser-vé ouvert ouver-te recou-vre-ment recou-vre-ments repré-sen-te}
\let\skpt\skippointrait
\let\epsilon\varepsilon
\let\ts\textstyle

\newenvironment{enumeratei}
{\bgroup\def\theenumi{\roman{enumi}}\def\theenumii{\arabic{enumii}}\begin{enumerate}}
{\end{enumerate}\egroup}
\newenvironment{enumeratea}
{\bgroup\def\theenumi{\alph{enumi}}\begin{enumerate}}
{\end{enumerate}\egroup}
\newcommand{\From}[1]{\mathchoice{\xleftarrow{\textstyle~#1~}}{}{}{}}
\newcommand{\isofrom}{\stackrel{\sim}{\longleftarrow}}

\newcommand{\alg}{\mathrm{alg}}
\newcommand{\CH}{\mathrm{CH}}

\usepackage{tikz}
\usepackage{tikz-cd}
\usepackage{csquotes}
\usepackage{adjustbox}
\usepackage{mathtools}

\newtheorem{thm}{Théorème}[section]
\newtheorem{prop}[thm]{Proposition}
\newtheorem{lemma}[thm]{Lemme}
\newtheorem{cor}[thm]{Corollaire}

\theoremstyle{definition}
\newtheorem{defin}[thm]{Définition}
\newtheorem{example}[thm]{Exemple}
\newtheorem{situation}[thm]{Situation}
\newtheorem{rmk}[thm]{Remarque}

\newcommand{\Q}{\mathbb Q}
\newcommand{\C}{\mathbb C}
\newcommand{\Z}{\mathbb Z}
\newcommand{\R}{\mathbb R}
\newcommand{\A}{\mathbb A}
\renewcommand{\P}{\mathbb P}
\newcommand{\Spec}{\operatorname{Spec}}
\newcommand{\Br}{\operatorname{Br}}
\newcommand{\Ker}{\operatorname{Ker}}

\newcommand{\mc}[1]{\mathcal{#1}}
\newcommand{\cl}{\overline}
\newcommand{\on}[1]{\operatorname{#1}}
\newcommand{\puiseux}{\mathbb{R}\{\!\{t\}\!\}}
\let\phi\varphi

\usetikzlibrary{decorations.markings}
\makeatletter
\tikzcdset{
open/.code = {\tikzcdset{hook, circled};},
closed/.code = {\tikzcdset{hook, slashed};},
open'/.code = {\tikzcdset{hook', circled};},
closed'/.code = {\tikzcdset{hook', slashed};},
circled/.code = {\tikzcdset{markwith = {\draw (0,0) circle (.375ex);}};},
slashed/.code = {\tikzcdset{markwith = {\draw[-] (-.4ex,-.4ex) -- (.4ex,.4ex);}};},
markwith/.code ={
\pgfutil@ifundefined%
{tikz@library@decorations.markings@loaded}%
{\pgfutil@packageerror{tikz-cd}{You need to say %
\string\usetikzlibrary{decorations.markings} to use arrows with markings}{}}{}%
\pgfkeysalso{/tikz/postaction = {
/tikz/decorate,
/tikz/decoration={markings, mark = at position 0.5 with {#1}}}
}
},
}
\makeatother


% End literal retained input author-preamble.tex

% Render-only compatibility: discard inert semicolons after style-setting calls.
% Drawing-path semicolons inside markwith remain intact. Native preamble/body unchanged.
\tikzcdset{
open/.code = {\tikzcdset{hook, circled}},
closed/.code = {\tikzcdset{hook, slashed}},
open'/.code = {\tikzcdset{hook', circled}},
closed'/.code = {\tikzcdset{hook', slashed}},
circled/.code = {\tikzcdset{markwith = {\draw (0,0) circle (.375ex);}}},
slashed/.code = {\tikzcdset{markwith = {\draw[-] (-.4ex,-.4ex) -- (.4ex,.4ex);}}},
}
\begin{document}
\begin{center}{\Large Stable item 2387}\end{center}
\paragraph{Source and attribution.}
Jean-Louis Colliot-Thélène, Alena Pirutka, Federico Scavia, \emph{Variétés réelles semi-algébriquement connexes non stablement rationnelles}, Journal de l’École polytechnique — Mathématiques 13 (2026), publisher version, DOI 10.5802/jep.343. \href{https://jep.centre-mersenne.org/articles/10.5802/jep.343/}{Publisher article}. Authors' text: \href{https://creativecommons.org/licenses/by/4.0/}{CC BY 4.0}.
\paragraph{Editorial problem boundary.}
Prove exactly the original French Theorem 3.6 below: the constructed smooth complete intersection of two quadrics in projective 5-space over the real Puiseux field has semi-algebraically connected real points and fails universal CH0 triviality, hence stable rationality. Retain the explicit real singular model, its four nonreal singularities, the complete norm-quotient specialization and finite-Puiseux descent, and the authored Nash-connectedness argument. The selected norm extension is the quadratic extension C/R used in this construction; no arbitrary-extension norm claim or analogue over the ordinary real field is selected. Cubic/higher-dimensional examples and later unramified-group calculations remain unselected context.
\paragraph{Difficulty (editorial): H3.}
A standard decomposition-of-diagonal specialization theorem does not by itself cover the singular special fiber used here. The authors pass to quotients by quadratic-extension norms, prove norm-surjectivity on the singular locus and its inverse image, and transport the nonconstant unramified Brauer obstruction through a generic-point DVR to all finite Puiseux extensions. They then construct a real two-quadrics model with only four nonreal singularities and a smooth deformation, while the complete Nash-triviality argument preserves semi-algebraic connectedness of the real locus.
\paragraph{Editorial transcription note.}
Original author language, mathematics and ordering are preserved. Publisher front matter is replaced by this independent article wrapper. Bibliography is recovered from matching cached publisher HTML, including its TeX formulas. No new proof or mathematical repair is supplied. Surrounding original results are retained for conservative context and are not extra selected corpus claims.
\clearpage

% Begin literal retained input author-body.tex
\section*{Introduction}

Soit $X$ une $\R$-variété projective lisse et géométriquement unirationnelle.
Pour $X$ de dimension $2$, on sait depuis Comessatti \cite{comessatti1912fondamenti} que si l'espace topologique $X(\R)$ est connexe (et en particulier non vide) alors $X$ est rationnelle sur $\R$. Ceci s'applique aux surfaces cubiques lisses dans $\P^3_{\R} $ et aux intersections lisses
de deux quadriques dans~$\P^4_{\R}$.
En dimension quelconque, que $X(\R)$ soit connexe
est une condition nécessaire pour que $X$ soit stablement rationnelle sur $\R$.
Ce n'est pas une condition suffisante de rationalité, déjà pour $X$
intersection lisse de deux quadriques dans $\P^5_\R$
\cite{hassett2021rationality}.
Mais la question de savoir si cette condition est suffisante pour assurer la rationalité stable est ouverte pour
certaines classes simples de variétés, par exemple les intersections lisses de deux quadriques dans $\P^5_\R$, ou~les hypersurfaces cubiques lisses de dimension au moins $3$.

On peut se poser la même question sur un corps réel clos $R$ arbitraire. On~dispose en effet d'une notion de connexité semi-algébrique qui étend la notion usuelle de connexité sur les réels; voir \cite{delfs1981semialgebraicI} et \cite[Def.\,2.4.2]{bochnak1998real}. Sur le corps \[\puiseux\coloneqq \bigcup_{n\geq1} \R(\!(t^{1/n})\!)\]
des séries de Puiseux réelles (qui est un corps réel clos avec $t$ infiniment petit positif) nous donnons des exemples
de variétés projectives lisses $X$ dont le lieu $X(\puiseux)$ des points réels est
semi-algébriquement connexe mais qui ne sont pas
$\CH_{0}$-triviales, et en particulier ne sont pas stablement rationnelles,
ni même rétractilement rationnelles, parmi les variétés des types suivants :

\begin{enumeratei}
\item\label{enum:i}
intersection lisse de deux quadriques dans $\P^5_{\puiseux}$ (théorème \ref{prop-intersection-quadriques}),
\item\label{enum:ii}
hypersurface cubique lisse dans $\P^4_{\puiseux}$ (théorème \ref{prop-cubique}),
\item\label{enum:iii}
solide lisse fibré en coniques sur $ \P^2_{\puiseux}$ et géométriquement rationnel (théorème \ref{prop-fibres-coniques}).
\end{enumeratei}
On ne sait pas s'il existe de tels exemples sur le corps $\R$ des réels.
Ces résultats sont à mettre en perspective avec plusieurs articles récents.

\begin{enumerate}
\item Sur tout corps $k$, des travaux de Hassett et Tschinkel \cite{hassett2021rationality} (pour $k=\R$) et de Benoist et Wittenberg \cite{benoist2019intermediate} (pour $k$ arbitraire) ont établi
la non-rationalité de toute intersection lisse $X$ de deux quadriques
dans $\P^5_{k}$, dès que $X$ ne contient
pas de droite définie sur $k$, ce qui sur $k=\R$ peut se produire avec $X(\R)$ connexe.
La méthode est une élaboration sur un corps arbitraire de celle de Clemens et
Griffiths pour établir la non-rationalité des hypersurfaces cubiques lisses dans $\P^4_{\C}$.
Olivier Wittenberg a récemment appliqué cette méthode pour établir sous des hypothèses générales la non-rationalité
de solides $X$ fibrés en surfaces quadriques sur la droite projective $\P^1_{\R}$, certains d'entre eux satisfaisant que $X(\R)$ est connexe. Cette méthode est puissante, mais elle est spécifique à la dimension $3$, et elle ne dit rien sur la rationalité stable.

\item Deux des auteurs \cite{colliot2024certaines} ont donné des exemples de solides $X$ fibrés en surfaces quadriques sur $\P^1_{\R}$
qui ne sont pas $\CH_{0}$-triviaux bien que $X(\R)$
soit connexe. L'invariant ici utilisé est le groupe de Brauer non ramifié.
Les fibrations sur $\P^1_{\R}$ ont des fibres géométriques singulières réductibles.
\item Deux des auteurs \cite{colliot2024certaines} ont étudié une classe de
solides $X$ fibrés en surfaces quadriques sur $\P^1_{\R}$, à fibres géométriques intègres, qui échappe aux deux techniques précédentes. Pour une sous-classe assez large de tels solides $X$ ils ont établi
que $X$ est $\CH_{0}$-triviale. La~question si cela vaut toujours pour toutes les $\R$-variétés de la classe est ouverte.
\item Sur le corps $\puiseux$, Benoist et le deuxième auteur \cite{benoist2026rationality} ont donné
des exemples de solides du type de \cite{colliot2024certaines}
dont l'espace des $\puiseux$-points est semi-algébriquement connexe et
qui ne sont pas $\CH_{0}$-triviaux. L'invariant utilisé
est un invariant en cohomologie non ramifiée de degré $3$ introduit dans \cite{colliot2024certaines}, reposant sur le travail \cite{colliot1993groupe}.
La non-nullité de cet invariant est détectée par
une méthode de spécialisation au-dessus des corps
$\R(\!(t^{1/n})\!)$ pour tous les $n\geq 1$.
\end{enumerate}

Dans le présent article, nous développons une nouvelle méthode, qui
donne des résultats pour plusieurs autres classes de variétés. Nous partons de versions singulières $Y/\R$ des solides $X$ des types \eqref{enum:i}, \eqref{enum:ii}, \eqref{enum:iii}, pour lesquelles
$Y(\R)$ est connexe, le lieu singulier est formé d'un nombre fini de points non réels,
et le groupe de Brauer non ramifié de $Y$ n'est pas réduit à celui de $\R$. De tels exemples ont été construits dans \cite{colliot2024certaines}.
Ensuite, nous déformons une telle variété $Y$ en une famille $\mc{X}$ sur $\A^1_\R=\Spec (\R[t])$ à fibre générique lisse du même type. En nous appuyant sur une version semi-algébrique du théorème d'Ehresmann due à Coste et Shiota \cite{coste1992nash}, nous montrons (paragraphe~\ref{ehres}) que l'ensemble des $\puiseux$-points de la $\puiseux$-variété projective lisse $X\coloneqq \mc{X}\times_{\R[t]}\puiseux$ est semi-algébriquement connexe. Une méthode alternative pour obtenir ce résultat (paragraphe~\ref{ehres}) repose sur le travail de Scheiderer \cite{scheiderer1994real}.
Par ailleurs nous montrons (paragraphe~\ref{parasp}) que la méthode de spécialisation, initiée par Voisin et généralisée par deux des auteurs \cite{colliot2016hypersurfaces}, s'adapte de façon souple à ce contexte
(sans résolution explicite des singularités)
et établit que la $\puiseux$-variété $X$ n'est pas
$\CH_0$-triviale et donc n'est pas stablement rationnelle. La~construction des exemples est détaillée au paragraphe~\ref{paraex}.

\medskip

La nouvelle méthode s'applique dans un cadre plus large (paragraphe~\ref{methode-2}),
où le lieu singulier de la variété auxiliaire $Y$ n'est pas
nécessairement fini, et où le groupe de Brauer non ramifié est remplacé par un groupe de cohomologie non ramifiée supérieur.

Nous obtenons ainsi (paragraphe~\ref{dimsup}) des exemples de $\puiseux$-variétés projectives lisses $X$ non $\CH_0$-triviales avec $X(\puiseux)$ semi-algébriquement connexe
parmi des variétés de dimension supérieure des types suivants :
\begin{enumeratei}\setcounter{enumi}{3}
\item\label{enum:iv} fibration en quadriques lisse de dimension relative $6$ sur $\P^1_{\puiseux}$ (théorème \ref{fibres-quadriques-6}),
\item\label{enum:v} intersection lisse de deux quadriques dans $\P^9_{\puiseux}$ (théorème \ref{intersection-quadriques-p9}).
\end{enumeratei}

Ici l'analyse de la fibre singulière utilisée dans la déformation est plus délicate, et~on a besoin d'une résolution des singularités explicite.

On comparera \eqref{enum:i} et \eqref{enum:v} avec la rationalité, au moins pour $R=\R$, des intersections lisses de deux quadriques dans $\P^4_R$ (Comessatti \cite{comessatti1912fondamenti}), dans $\P^6_{R}$ et dans
un certain nombre de cas dans $\P^{2n}_R$ (Hassett, Kollár, Tschinkel \cite{hassett2022rationality}) lorsque leur lieu réel est semi-algébriquement connexe.

On conclut cet article avec des calculs (paragraphe~\ref{calcul}) des groupes de cohomologie non ramifiée en bas degré pour les variétés projectives et lisses ici considérées. Dans beaucoup de cas, ces groupes sont réduits à la cohomologie du corps de base, et ce même
après toute extension du corps de base,
si bien qu'ils ne permettent pas de détecter la non $\CH_0$-trivialité.

\subsection*{Notations et rappels}
Un corps $F$ est dit {\it formellement réel} s'il satisfait l'une des propriétés équivalentes suivantes (\cf \cite[Th.\,1.1.8]{bochnak1998real}) :
\begin{enumeratei}
\item
le corps $F$ peut être ordonné;
\item
l'élément $-1$ n'est pas une somme de carrés dans $F$;
\item
pour tous $x_1,\ldots x_n\in F$, si $\sum x_i^2=0$, alors $x_i=0$, $i=1,\ldots n$.
\end{enumeratei}

Un corps formellement réel $R$ est dit {\it réel clos} s'il n'admet pas d'extension algébrique réelle non triviale. Il est équivalent de dire (\cf \cite[Th.\,1.2.2]{bochnak1998real}) que le corps $R[i]\coloneqq R[x]/(x^2+1)$ est algébriquement clos.
Par exemple, les corps suivants sont des corps réels clos: le corps des nombres réels $\mathbb R$, le corps $\mathbb R_{\alg}$ des nombres réels qui sont algébriques sur $\mathbb Q$, et le corps des séries de Puiseux à coefficients réels $\puiseux$ (voir \cite[Ex.\,1.2.3]{bochnak1998real}).

\medskip

Soit $k$ un corps.
Une $k$-variété est un $k$-schéma séparé de type fini. Une $k$-variété intègre
est dite {\it $k$-rationnelle} si elle est $k$-birationnelle à un espace projectif $\P^n_{k}$.
Une $k$-variété intègre $X$ est dite {\it stablement $k$-rationnelle} s'il existe des espaces projectifs
$\P^n_{k}$ et $\P^m_{k}$ tels que $X \times_{k}\P^n_{k}$ soit $k$-birationnel à $\P^m_{k}$.
Une $k$-variété intègre $X$ est dite {\it rétractilement rationnelle}
s'il existe des ouverts de Zariski non vides $U \subset X$ et $V \subset \P^m_{k}$ ($m$ convenable),
et des $k$-morphismes $ f \colon U \to V$ et $g \colon V \to U$ tels que le composé $g \circ f$ soit l'identité de $U$.
Une $k$-variété intègre stablement $k$-rationnelle est rétractilement rationnelle.

\medskip

Pour $X$ une $k$-variété, on note $\CH_{0}(X)$ le groupe de Chow des zéro-cycles modulo l'équivalence rationnelle.
Pour $X$ une $k$-variété propre, l'application qui à un point fermé $P$
associe le degré $[k(P):k]$ du corps résiduel $k(P)$ s'étend en un homomorphisme
$ \mathrm{deg}_{k} \colon \CH_{0}(X) \to \Z$. On~note $A_{0}(X)$ le noyau de cet homomorphisme.

On dit qu'une $k$-variété propre géométriquement intègre $X$ est (universellement) $\CH_{0}$-triviale si pour tout corps $F$ contenant $k$
la flèche $ \mathrm{deg}_{F} \colon \CH_{0}(X_{F}) \to \Z$ est un isomorphisme
(voir \cite{auel2017universal, colliot2016hypersurfaces}).
Sous l'hypothèse que la $k$-variété $X$, de corps des fonctions $k(X)$, est lisse et
possède un zéro-cycle de degré $1$, ceci est équivalent à l'énoncé : $A_{0}(X_{k(X)})=0$ \cite[Lem.\,1.3]{auel2017universal}.

Si une $k$-variété géométriquement intègre propre et lisse est stablement rationnelle, ou~plus généralement rétractilement rationnelle, alors elle est $\CH_{0}$-triviale.

Soit $C$ un corps algébriquement clos de caractéristique zéro.
Rappelons que toute intersection complète lisse de deux quadriques dans $\P^n_C$ avec $n \geq 4$ est rationnelle,
et qu'il en est de même de toute variété projective et lisse de dimension au moins~2 fibrée en quadriques sur la droite projective $\P^1_C$. Il existe des solides fibrés en coniques sur le plan projectif qui ne sont pas rétractilement rationnels (Artin-Mumford).
Une hypersurface cubique lisse dans $\P^4_C$ n'est pas rationnelle (Clemens-Griffiths).

Soient $k$ un corps et $M$ un module galoisien discret de torsion, première à la caractéristique de $k$.
Soient
$X$ une $k$-variété intègre et $k(X)$
son corps des fonctions. On~note \[
H^{i}_{\on{nr}}(k(X)/k, M) \subset H^{i}(k(X),M)
\]
le sous-groupe du groupe de cohomologie galoisienne
formé des éléments qui sont non ramifiés par rapport à toute valuation discrète sur $k(X)$ triviale sur $k$; voir \hbox{\cite[\S 4]{colliot1995birational}}, \cite{rost1996chow} et \cite{merkurjev2008unramified}.
Si $\on{car}(k)=0$ et $X/k$ est projective et lisse, on a
\[\on{Br}(X)=\on{Br}_{\on{nr}}(k(X)/k)=H^2_{\on{nr}}(k(X)/k,\Q/\Z(1)),\]
où $\on{Br}(X)=H^2_{\text{ét}}(X,\mathbb{G}_m)$ est le groupe de Brauer de $X$, où $\on{Br}_{\on{nr}}(k(X)/k)$ est le groupe de Brauer non ramifié de $k(X)$ sur $k$ et où, pour tout $j\in \Z$, on note $\Q/\Z(j)\coloneqq \varinjlim_n \mu_n^{\otimes j}$.

On utilisera dans cet article des résultats de la théorie algébrique des formes quadratiques \cite{lam2005introduction, kahn2008formes}. En particulier, on utilisera les propriétés des formes de Pfister (voir \cite[Ch.\,X]{lam2005introduction} et \cite[Ch.\,2]{kahn2008formes}) et des formes d'Albert (voir \cite[(4.7) p.\,69, Albert's Th.\,4.8]{lam2005introduction} et \cite[Déf.\,5.7.7, Lem.\,8.1.4]{kahn2008formes}).

\subsubsection*{Remerciements}
Les auteurs remercient Olivier Wittenberg pour avoir suggéré la deuxième preuve du théorème \ref{critere-connexe}.

\section{Un critère de connexité semi-algébrique}\label{ehres}

Dans cette section, nous donnons un critère de connexité semi-algébrique, au sens de \cite[Def.\,2.4.2]{bochnak1998real}, pour la fibre générique, au-dessus du corps des séries de Puiseux $\puiseux$, d'un morphisme projectif de $\R$-variétés $f\colon X\to U$, où $U\subset \A^1_\R=\Spec \R[t]$ est un sous-schéma ouvert contenant $0$.

\begin{thm}\label{critere-connexe}
Soit $\A^1_\R=\Spec \R[t]$ et soit $\eta\colon \Spec(\puiseux)\to \A^1_\R$ le morphisme correspondant à l'inclusion $\R[t]\subset \puiseux$. Soit $U\subset \A^1_\R$ un voisinage Zariski de $0$, soit $f\colon X\to U$ un morphisme projectif fidèlement plat de $\R$-variétés, à fibre générique lisse, soit $X_0\coloneqq X|_{t=0}$ et soit $X_\eta\coloneqq X\times_{U,\eta}\Spec(\puiseux)$. Supposons que $X_0(\R)$ est non vide et contenu dans le lieu lisse de $X_0$. Le nombre de composantes connexes de $X_0(\R)$ coïncide avec le nombre de composantes semi-algébriquement connexes de $X_\eta(\puiseux)$.
\end{thm}

On donnera deux preuves de ce résultat. La~première démonstration est analytique et utilise un théorème fondamental de Nash-trivialité des fibrations dû à Coste et Shiota \cite{coste1992nash}. La~deuxième démonstration utilise la variante réelle de la théorie de SGA 4, développée par Scheiderer \cite{scheiderer1994real}.

\subsection{Première preuve du théorème \ref{critere-connexe}}

Nous introduisons maintenant les notations nécessaires pour énoncer le théorème de Nash-trivialité des fibrations
de Coste et Shiota \cite{coste1992nash}; nos références de base sont l'article \cite{delfs1981semialgebraicII} de Delfs et Knebusch et le livre \cite{bochnak1998real} de Bochnak, Coste et Roy.

Soit $R$ un corps réel clos et soit $X$ une $R$-variété. D'après \cite[Ex.\,2 p.\,182]{delfs1981semialgebraicII} l'ensemble $X(R)$ admet une structure naturelle d'espace semi-algébrique sur $R$, au sens de \cite[Def.\,3 p.\,182]{delfs1981semialgebraicII} (il s'agit donc d'espaces annelés sur $R$ au sens de \cite[Def.\,2 p.\,182]{delfs1981semialgebraicII}). Si $X$ est quasi-projective, $X(R)$ est même un espace semi-algébrique affine sur $R$; voir \cite[Rem.\ p.\,182]{delfs1981semialgebraicII}. Pour tout morphisme de $R$-variétés $f\colon X\to Y$, l'application induite $f(R)\colon X(R)\to Y(R)$ est semi-algébrique au sens de \cite[p.\,184]{delfs1981semialgebraicII}. Si le morphisme $f$ est propre, d'après \cite[Th.\,9.6]{delfs1981semialgebraicII} l'application $f(R)$ est propre au sens de \cite[Def.\ p.\,192]{delfs1981semialgebraicII}, et en particulier elle est fermée.

Soit $X$ une $R$-variété quasi-projective. Alors $X(R)$ est un espace semi-algébrique affine sur $R$. On~peut donc exhiber $X(R)$ comme un sous-ensemble algébrique \hbox{$X(R)\subset R^n$} au sens de \cite[Def.\,2.1.1]{bochnak1998real}, pour $n\geq 1$ convenable. Supposons de plus que $X(R)$ est contenu dans le lieu lisse de $X$. Alors, d'après \cite[Prop.\,3.3.11]{bochnak1998real}, $X(R)$ est donc une variété de Nash affine au sens de \cite[p.\,351]{coste1992nash}, ou~de façon équivalente une sous-variété de Nash de $R^n$ au sens de \cite[Def.\,2.9.9]{bochnak1998real}. Si $Y$ est une autre $R$-variété quasi-projective telle que $Y(R)$ soit contenu dans le lieu lisse de~$Y$ et $f\colon X\to Y$ est un $R$-morphisme tel que $X(R)$ soit contenu dans le lieu lisse de~$f$, l'application semi-algébrique induite $f(R)$ est une submersion de Nash, c'est-à-dire une application de Nash au sens de \cite[Def.\,2.9.9]{bochnak1998real} dont la différentielle est surjective en tout point de $X(R)$. En particulier, $f(R)$ est ouverte. On~rappelle que, par définition, toute application de Nash $p\colon M\to N$ entre variétés de Nash affines est en particulier une application semi-algébrique, et que, par définition, $p$ est un difféomorphisme de Nash si elle est bijective et $p^{-1}$ est aussi une application de Nash.

\begin{thm}[Coste-Shiota]\label{coste-shiota-thom}
Soit $R$ un corps réel clos, soit $M$ une variété de Nash affine, et soit $p\colon M\to R$ une submersion propre de Nash. Alors il existe un difféomorphisme de Nash $h\colon p^{-1}(0)\times R\to M$ tel que $p\circ h$ soit la projection sur $R$.
\end{thm}

\begin{proof}
On peut supposer $M$ non vide. Comme toute submersion est ouverte et toute application propre est fermée, l'image de $p$ est ouverte, fermée et non vide, et donc $p$ est surjective. D'après \cite[Th.\,A, Th.\,2.4(iii')]{coste1992nash}, il~existe alors une variété de Nash affine $F$ et un difféomorphisme de Nash $h'\colon F\times R\to M$ tel que $p\circ h'$ soit la projection sur $R$. Par restriction à $0\in R$, on déduit un difféomorphisme de Nash $p_0\colon F\to p^{-1}(0)$. Il suffit de poser $h\coloneqq h'\circ (p_0^{-1}\times \on{id}_R)$ pour conclure. (Par convention, toutes les variétés de Nash considérées dans \cite{coste1992nash} sont affines; voir \cite[p.\,351]{coste1992nash}.)
\end{proof}

\begin{proof}[Première démonstration du théorème \ref{critere-connexe}]
Soit $R$ le corps $\puiseux$. D'après \cite[Th.\,2.4.5]{bochnak1987geometrie}, les composantes connexes et les composantes semi-algébriquement connexes de $X_0(\R)$ coïncident. Par \cite[Prop.\,5.3.6]{bochnak1998real}, $X_0(\R)$ et $X_0(R)$ ont le même nombre de composantes semi-algébriquement connexes. Il suffit donc de montrer que les semi-algébriques $X_0(R)$ et $X_\eta(R)$ ont le même nombre de composantes semi-algébriquement connexes. On~prouvera que $X_0(R)$ et $X_\eta(R)$ sont Nash-difféomorphes entre eux: le théorème en découle parce qu'un difféomorphisme de Nash est ouvert et fermé et l'image d'un espace semi-algébriquement connexe par une application semi-algébrique est encore semi-algébriquement connexe.

Comme la fibre générique du morphisme $f$ est lisse, passant à un ouvert de Zariski $U\subset \A^1_{\R}$ contenant $0$ plus petit si nécessaire, on peut supposer que $f$ est lisse au-dessus de $U\setminus\{0\}$. Soit $V\subset X$ le plus grand sous-schéma ouvert de $X$ tel que la restriction de $f$ à $V$ soit lisse. Comme le morphisme $f$ est plat, l'hypothèse que~$X_0(\R)$ est contenu dans le lieu lisse de $X_0$ entraîne que $X_0(\R)$ est contenu dans $V$. On~a donc $V(\R)=X(\R)$, et le théorème d'homomorphisme d'Artin-Lang \cite[Th.\,4.1.2]{bochnak1998real} entraîne alors $V(R)=X(R)$. En particulier, $X(R)$ est une variété de Nash affine sur $R$.
De plus, l'application $f(R)$ est une submersion de Nash propre : c'est une submersion car $V(R)=X(R)$ et elle est propre car le morphisme $f_R$ est propre.

Comme l'ouvert $U\subset \A^1_{\R}$ contient $0$, il~existe $\epsilon\in \R_{>0}$ tel que l'intervalle ouvert $(-\epsilon,\epsilon)\subset \R$ soit contenu dans $U(\R)$. Si $Y\subset \A^1_{\R}$ est le complémentaire de $U$, l'appli\-cation $Y(\R)\to Y(R)$ est bijective (c'est un cas très facile du théorème d'homomorphisme d'Artin-Lang). On~en déduit que l'élément $t\in R=\A^1_R(R)$ appartient à~$U(R)$ et que l'intervalle ouvert $(-\epsilon,\epsilon)\subset R$ est contenu dans $U(R)$. Le morphisme $\eta\colon \operatorname{Spec}(R) \to U$ se factorise par le morphisme $\operatorname{Spec}(R) \to U_R$ correspondant à l'élément $t\in U(R)$. Comme $\epsilon\in \R_{>0}$ et $t$ est infinitésimal, on a $-\epsilon < t<\epsilon$. L'intervalle $(-\epsilon,\epsilon)\subset R$ étant Nash-difféomorphe à $R$, le théorème \ref{coste-shiota-thom} entraîne alors que $X_0(R)$ et $X_\eta(R)$ sont Nash-difféomorphes entre eux, comme voulu.
\end{proof}

\subsection{Deuxième preuve du théorème \ref{critere-connexe}}
Le théorème \ref{critere-connexe} est un cas particulier du
théorème suivant, qui est une légère extension de \cite[Cor.\,17.20(a)]{scheiderer1994real}. Pour tout schéma $X$, on note par $X_r$ son spectre réel, avec la topologie engendrée par les domaines de positivité \cite[(0.4.1)--(0.4.2)]{scheiderer1994real}, et on écrit $X_{ret}$ pour son site étale réel \cite[Def.\,1.2.1]{scheiderer1994real}, qui est la catégorie $\mathrm{Et}/X$ des $X$-schémas étales avec les recouvrements donnés par les morphismes réels surjectifs au sens de \cite[Def.\,(1.1)]{scheiderer1994real}, c'est-à-dire les collections de morphismes $\{f_i\colon U_i\to U\}$ dans $\mathrm{Et}/X$ telles que $U_r$ soit la réunion des $f_{ir}(U_{ir})$. D'après \cite[Th.\,1.3]{scheiderer1994real}, les topoï associés à $X_r$ et $X_{ret}$ sont naturellement équivalents entre eux. Par définition, un faisceau abélien $\mc{F}$ sur $X_r$ est dit constructible si $X_r$ est la réunion d'un nombre fini de $K_i\subset X_r$ constructibles tels que $\mc{F}|_{K_i}$ soit localement constant avec fibres de type fini; voir \cite[Def.\,(A.3)]{scheiderer1994real}.

\begin{thm}\label{loc-const}
Soit $R$ un corps réel clos et soit $f\colon X\to U$ un morphisme propre de $R$-variétés. Supposons ce morphisme lisse en tout $x\in X(R)$. Soit $\mc{F}$ un faisceau abélien constructible localement constant sur $X_r$. Alors, pour tout $n\geq 0$, le faisceau $R^n f_{r*} \mc{F}$ est constructible et localement constant sur $U_r$.
\end{thm}

\begin{proof}
Soient $\eta$ et $\xi$ deux points de $U$ tels que $\xi$ soit une spécialisation de $\eta$ et soit $\mc{G}$ un faisceau sur $U_r$. On~note par $\mc{G}_{\xi}$ et $\mc{G}_{\eta}$ les fibres de $\mc{G}$ en~$\xi$ et~$\eta$, respectivement. Soit $V\subset U_r$ un voisinage ouvert de $\xi$. Alors $V$ contient $\eta$ et on dispose donc d'une application naturelle $\mc{G}(V)\to \mc{G}_\eta$. Par passage à la limite inductive en $V$, on obtient une flèche canonique
\[\on{cosp}_{\eta\leadsto\xi}\colon \mc{G}_{\xi} \to \mc{G}_{\eta}\]
qui est dite application de cospécialisation; voir \cite[(5) p.\,203]{scheiderer1994real} et \cite[Ch.\,VIII, 7.7]{sga4II}. Par construction, pour tout voisinage ouvert $V$ de $\xi$ et tout voisinage ouvert $W$ de $\eta$ contenu dans $V$, on a un diagramme commutatif
\begin{equation}
\begin{tikzcd}[column sep = large,row sep = large]\label{cosp-square}
\mc{G}(V) \arrow[r] \arrow[d] & \mc{G}(W) \arrow[d] \\
\mc{G}_\xi \arrow[r,"\ts\on{cosp}_{\eta\leadsto\xi}"] & \mc{G}_\eta.
\end{tikzcd}
\end{equation}

On revient à la preuve du théorème \ref{loc-const}. Si le morphisme propre $f$ était lisse, l'énoncé suivrait de \cite[Cor.\,17.20(a)]{scheiderer1994real}. On~adapte le même argument au cas général. Par \cite[Th.\,17.7]{scheiderer1994real}, pour tout $n\geq 0$ le faisceau $R^n f_{r*} \mc{F}$ est constructible. Par conséquent, il~suffit de montrer que pour toute spécialisation $\eta \leadsto \xi$ dans $U_r$, les applications de cospécialisation\enlargethispage{-\baselineskip}
\[\on{cosp}_{\eta\leadsto\xi}\colon (R^n f_{r*} \mc{F})_{\xi} \to (R^n f_{r*} \mc{F})_{\eta}\] sont des isomorphismes. On~peut remplacer $U$ par la clôture schématique de $\{\eta\}$ et donc se réduire au cas où $U$ est intègre avec point générique $\eta$.

Soit $B$ l'enveloppe convexe de l'anneau local $O_{U,\xi}$ dans une clôture réelle $R'$ du corps des fractions de $U$; donc $B$ est un anneau de valuation réel clos au sens de \cite[(1.5)]{scheiderer1994real}. L'espace topologique $\on{Spec}(B)$ est ordonné linéairement par la relation de spécialisation; on montre qu'il est fini. Soit $p_0\subsetneq p_1\subsetneq\dots p_m$ une chaîne d’idéaux premiers de $B$ de longueur $m$, soit $x_i\in p_i\setminus p_{i-1}$ pour tout $i=1,\dots, m$ et soit $C\subset B$ la sous-$R$-algèbre de $B$ engendrée par les $x_i$. Les idéaux $p_i\cap C$ forment une chaîne d’idéaux premiers de $C$ de longueur $m$ et donc $m\leq \on{dim}(C)$. Comme $C$ est de type fini sur $R$, par le théorème de la dimension $\on{dim}(C)=\on{tr}_R(\on{Frac}(C))\leq \on{tr}_R(R')$, où $\on{tr}$ est le degré de transcendance. On~conclut que $m\leq \on{tr}_R(R')$ et donc que l'ensemble sous-jacent à $\on{Spec}(B)$ est fini. En particulier, le point générique de $\on{Spec}(B)$ est ouvert.

Le morphisme composé $\Spec(B) \to \on{Spec}(O_{U,\xi})\to U$ envoie le point générique de $\on{Spec}(B)$ vers $\eta$ et le point fermé de $\on{Spec}(B)$ vers $\xi$. Comme le morphisme $f$ est propre et de présentation finie, par le théorème de changement de base propre pour la topologie réelle étale \cite[Th.\,16.2(a), $t=ret$]{scheiderer1994real} on peut remplacer $U$ par $\Spec(B)$. Donc $U = \Spec(B)$, où $B$ est un anneau de valuation réel clos, $\eta$ est le point générique de $U$ et $\xi$ est le point fermé de $U$. On~observe que $U$ (et donc $X$) n'est plus une $R$\nobreakdash-variété en général.

Soit $g\colon \eta \to U$ l'immersion ouverte donnée par l'inclusion du point générique de $U$ et soit $X_\eta$ la fibre générique de $f$, de sorte que l'on a un carré cartésien
\[
\begin{tikzcd}[column sep = large,row sep = large]
X_\eta \arrow[r,"\ts h"] \arrow[d,swap,"\ts f_\eta"] & X \arrow[d,"\ts f"] \\
\eta \arrow[r,"\ts g"] & U
\end{tikzcd}
\]
où les flèches horizontales sont des immersions ouvertes. Le diagramme commutatif~\eqref{cosp-square} pour $\mc{G}=R^nf_*\mc{F}$, $V=U$ et $W=\{\eta\}$ prend la forme\vspace*{-5pt}
\[
\begin{tikzcd}[column sep = large]
H^n(X_r, \mc{F}) \arrow[r,"\ts h_r^*"] \arrow[d,"\ts \wr"] & H^n(X_{\eta r}, h_r^* \mc{F}) \arrow[d,"\ts \wr"] \\
(R^n f_{r*} \mc{F})_{\xi} \arrow[r,"\ts \on{cosp}_{\eta\leadsto\xi}"] & (R^n f_{r*} \mc{F})_{\eta}
\end{tikzcd}
\]
Ici la flèche du haut coïncide avec l'application induite par l'immersion ouverte $h_r$. La~flèche verticale de gauche est un isomorphisme car $U_r$ (\resp $\{\eta\}$) est le plus petit ouvert de $U_r$ contenant $\xi$ (\resp $\eta$). (On rappelle que d'après \cite[(1.5)]{scheiderer1994real}, l'espace topologique sous-jacent au schéma $U$ est homéomorphe à $U_r$.)

Il suffit donc de montrer que l'application $h_r^*\colon H^n(X_r, \mc{F})\to H^n(X_{\eta r}, h_r^* \mc{F})$ est un isomorphisme. Cette application est un homomorphisme de coin dans la suite spectrale de Leray pour le morphisme $h_r\colon X_{\eta r}\to X_r$ et le faisceau $h_r^*\mc{F}$ sur $X_{\eta r}$:\vspace*{-3pt}
\[E_2^{pq}\coloneqq H^p(X_r,R^qh_{r*}h_r^{*}\mc{F})\Longrightarrow H^{p+q}(X_{\eta r},h_r^*\mc{F}).\]
Il suffit alors de prouver que la flèche canonique
$\mc{F} \to h_{r*}h_r^* \mc{F}$ est un isomorphisme et que $R^q h_{r*} h_r^* \mc{F} = 0$ pour tout $q\geq 1$.
Ces deux assertions sont locales sur $X_r$ par rapport à sa topologie réelle.
Soit $(P,\alpha)\in X_r$ un point, c'est-à-dire que l'on a un point schématique $P\in X$ et un ordre $\alpha$ sur le corps résiduel de $P$. Les hypothèses entraînent que le morphisme $f$ est lisse en $P$. Pour montrer que les deux assertions sont satisfaites au voisinage de $P$, on peut alors remplacer $X$ par l'ouvert maximal de lissité de $f$. Donc, renonçant ainsi à la propreté de $f$, on peut supposer que le morphisme $f$ est lisse.

Remplaçant $X$ par un recouvrement $X'\to X$ dans $X_{ret}$ (c'est-à-dire, un morphisme étale tel que le morphisme induit $X'_r\to X_r$ soit surjectif) tel que $\mc{F}|_{X'}$ soit constant, on se réduit au cas où $\mc{F}$ est constant de fibre un groupe abélien de type fini $M$. Alors\vspace*{-3pt}
\[ R^q h_{r*} h_r^* \mc{F} = R^q h_{r*} M = R^q h_{r*}f_{\eta r}^* M \isofrom f_r^* R^q g_{r*} M, \]
le dernier isomorphisme étant donné par le théorème de changement de base lisse pour la topologie $t=ret$; voir
\cite[Th.\,16.11]{scheiderer1994real}. Pour appliquer ce théorème, il~faut rappeler que, pour tout schéma $S$, tout faisceau abélien sur $S_{ret}$ est admissible \cite[Def.\,16.4.1(a)]{scheiderer1994real}. Comme $\eta$ est le spectre d'un corps réel clos, le foncteur $g_{r*}$ est exact: si $\mc{F}_1\to \mc{F}_2$ est un morphisme surjectif de faisceaux sur $\eta$, alors comme $\eta$ n'admet pas de recouvrement non trivial dans $(\eta)_{ret}$, le morphisme $\mc{F}_1(\eta)\to \mc{F}_2(\eta)$ des sections globales est surjectif, c'est-à-dire, $\mc{F}_1\to \mc{F}_2$ est surjectif en tant que morphisme de préfaisceaux, et donc il en est de même pour le morphisme $g_{r*}\mc{F}_1\to g_{r*} \mc{F}_2$. Puisque $g_{r*}M$ est le faisceau constant $M$ sur $U_r$, ceci achève la démonstration.
\end{proof}

\begin{proof}[Deuxième démonstration du théorème \ref{critere-connexe}]
Les composantes connexes de $X_0(\R)$ coïncident avec ses composantes semi-algébriquement connexes; voir \cite[Th.\,2.4.5]{bochnak1987geometrie}. Soit $m_0$ (\resp $m_\eta$) le nombre de composantes connexes semi-algébriques de $X_0$ (\resp $X_\eta$). Le théorème \ref{loc-const} et le théorème de changement de base propre \cite[Th.\,16.2 ($t=ret$)]{scheiderer1994real}, appliqués au faisceau constant $\mc{F}=\Z_{X_r}$ et $n=0$, donnent un isomorphisme entre $H^0(X_{0r},\Z)\cong \Z^{m_0}$ et $H^0(X_{\eta r},\Z)\cong\Z^{m_\eta}$, ce qui entraîne $m_0=m_\eta$.
\end{proof}

\subsection{Invariance birationnelle du nombre de composantes semi-algébriques connexes} Pour calculer le nombre de composantes connexes de certaines variétés singulières, le lemme suivant nous sera utile.

\begin{lemma}\label{connexe-invariant}
Soit $R$ un corps réel clos. Le nombre de composantes semi-algébri\-quement connexes de $X(R)$ est un invariant birationnel des $R$-variétés projectives intègres $X$ telles que $X(R)$ soit contenu dans le lieu lisse de $X$.
\end{lemma}

\begin{proof}
L'invariance birationnelle du nombre de composantes semi-algé\-briquement connexes des $R$-points des $R$-variétés projectives lisses est un fait classique \cite[\S 13]{delfs1981semialgebraicII}, \cite[Th.\,3.4.12]{bochnak1998real}. Il suffit donc de montrer que, pour toute $R$-variété projective intègre $X$ telle que $X(R)$ soit contenu dans le lieu lisse de $X$, il~existe une variété projective lisse $X'$, birationnelle à $X$, telle que $X(R)$ et $X'(R)$ aient le même nombre de composantes semi-algébriquement connexes. D'après Hironaka, il~existe une variété projective lisse $X'$ et un morphisme projectif birationnel $X'\to X$ qui est un isomorphisme au-dessus du lieu lisse de $X$. Comme le lieu lisse de $X$ contient $X(R)$, on déduit que $X(R)$ et $X'(R)$ sont semi-algébriquement isomorphes, et donc ils ont le même nombre de composantes semi-algébriquement connexes, comme voulu.
\end{proof}

\begin{rmk}\label{fibres-connexes}
Soient $R$ un corps réel clos et $f\colon X\to Y$ un morphisme fermé ou ouvert d'espaces semi-algébriques sur $R$. Supposons que pour tout $y\in Y$ la fibre $f^{-1}(y)$ est semi-algébriquement connexe. Si $Y$ est semi-algébriquement connexe, alors~$X$ l'est aussi d'après \cite[Exer.\,4.4.1(a)]{scheiderer2024course}. Plus généralement, si $Y_1,\dots,Y_r$ sont les composantes semi-algébriquement connexes de $Y$ et, pour tout $i=1,\dots,r$, on pose $X_i\coloneqq f^{-1}(Y_i)$, alors $X_1,\dots,X_r$ sont les composantes semi-algébriquement connexes de $X$.
\end{rmk}

\section{La méthode de spécialisation revisitée, I}
\label{parasp}

\subsection{Sur un corps quelconque}

La méthode suivie est une variante de la méthode de spécialisation sur un corps quelconque, sous la forme
développée dans \cite[\S 1]{colliot2016hypersurfaces}.

\begin{situation}\label{situation-phi}
Soit $L/k$ une extension de corps de caractéristique zéro, soit $Y$ une $k$-variété propre et géométriquement intègre, soit $S\subset Y$ le lieu singulier de $Y$ (avec structure réduite), soit $W\coloneqq Y\setminus S$ l'ouvert complémentaire de $S$ et soit \hbox{$b\in W(k)$}. Soit encore $Z$ une $k$-variété propre lisse géométriquement connexe, soit \hbox{$p\colon Z \to Y$} un $k$-morphisme propre birationnel tel que la restriction $p^{-1}(W) \to W$ soit un isomorphisme, et posons $T\coloneqq p^{-1}(S) \subset Z$. (Comme le corps $k$ est de caractéristique zéro, un tel morphisme $p$ existe d'après le théorème d'Hironaka.) Notons $E\coloneqq k(Y)=k(Z)$ et $E'\coloneqq L(Y)=L(Z)\cong E\otimes_kL$. Soit
\begin{equation}\label{application-phi}\Phi\colon \CH_{0}(Z_{E})/ N_{E'/E}(\CH_{0}(Z_{E'})) \to \CH_{0}(Y_{E})/N_{E'/E}(\CH_{0}(Y_{E'}))
\end{equation}
l'application induite par le morphisme propre $p$.
\end{situation}
La situation \ref{situation-phi} est résumée par le diagramme commutatif suivant, dans lequel les carrés sont cartésiens.
\[
\begin{tikzcd}
T\arrow[r, closed] \arrow[d] & Z \arrow[d,"\ts p"] & \arrow[l,open'] p^{-1}(W) \arrow[d,"\ts \wr"] \\
S \arrow[r,closed] & Y & \arrow[l,open'] W
\end{tikzcd}
\]

\begin{prop}\label{prop1}
On se place dans la situation \ref{situation-phi}. Soit $G_k$ le groupe de Galois absolu de $k$, soit $M$ un $G_k$-module discret de torsion et soit $i\geq 0$ un entier. Supposons:

\begin{enumeratei}
\item\label{prop1i}
L'application $\Phi$ de \eqref{application-phi} est injective.
\item\label{prop1ii}
L'application \[\Ker[H^i(k,M) \to H^i(L,M)] \to \Ker [H^i_{\on{nr}}(k(Y)/k,M) \to H^i_{\on{nr}}(L(Y)/L,M)]\] n'est pas
surjective.
\end{enumeratei}
Alors la différence entre le point générique de $Y$
et le point $b_{k(Y)}$ définit une classe non nulle dans $A_{0}(Y_{k(Y)})
\subset \CH_{0}(Y_{k(Y)})$.
\end{prop}

\begin{proof}
L'application $\Phi$ envoie la classe du point générique de
$Z$ (vu comme $E$-point de $Z_E$) sur celle du point générique de $Y$ (vu comme $E$-point de $Y_E$). D'après l'hypothèse \eqref{prop1ii}, il~existe $\alpha \in \Ker [H^i_{\on{nr}}(E/k,M) \to H^i_{\on{nr}}(E'/L,M)]$ qui ne vient pas de $\Ker[H^i(k,M) \to H^i(L,M)]$. Soit $\alpha|_b\in H^i(k,M)$ la spécialisation de $\alpha$ en $b$. Comme $\alpha_{L(Y)}=0$ et la spécialisation commute avec la restriction, la classe $\alpha|_b$ appartient à $\Ker[H^i(k,M) \to H^i(L,M)]$. En remplaçant $\alpha$ par $\alpha-\alpha|_b$, on peut supposer que $\alpha$ est non nul, mais nul en $b$ et donc au point $b'$ image inverse de $b$ par $p\colon Z \to Y$. D'après Merkurjev \cite[Cor.\,2.9]{merkurjev2008unramified} (\cf \cite[Cor.\,5.2]{schreieder2021unramified}), on dispose
de l'accouplement
\[\CH_{0}(Z_{E})/ N_{E'/E}(\CH_{0}(Z_{E'})) \times \Ker [H^i_{\on{nr}}(E/k,M) \to H^i_{\on{nr}}(E'/L,M)] \to H^i(E,M)\]
qui associe au point générique de $Z$ (respectivement, au point $b'_{E}\coloneqq b'\times_k E$) la classe $\alpha \in H^{i}(E,M)$ (respectivement, zéro, par le choix de la classe $\alpha$). Ainsi la différence entre le point générique de $Z$ et le point constant $b'_E$
définit une classe non triviale dans $\CH_{0}(Z_{E})/ N_{E'/E}(\CH_{0}(Z_{E'}))$. D'après l'hypothèse \eqref{prop1i}, on en déduit que
la différence entre le point générique de $Y$ (vu comme $E$-point de $Y_E$) et le point $b_{E}\coloneqq b\times_{k}E$
définit une classe de degré zéro qui est non nulle dans $\CH_{0}(Y_{E})/ N_{E'/E}(\CH_{0}(Y_{E'}))$,
et a fortiori dans $\CH_{0}(Y_{E})$.
\end{proof}

Le théorème suivant est dans l'esprit de \cite[Th.\,1.14]{colliot2016hypersurfaces}.

\begin{thm}\label{1.14K}
On se place dans la situation \ref{situation-phi}. Soit $A$ un anneau de valuation discrète de corps résiduel $k$, soit $K$ son corps des fractions, soit $\pi \colon { \mathcal{X}} \to \Spec(A)$ un $A$\nobreakdash-schéma projectif, fidèlement plat, à fibres géométriquement intègres, tel que \hbox{$\mc{X}\times_Ak\cong Y$}, soit $X$ la fibre générique de $\pi$ et soit $\sigma$ une section de $\pi$ qui induit un $k$-point lisse $b$ sur $Y$. Soit $M$ un $G_k$-module discret, de torsion, et soit $i\geq 0$ un entier. Supposons:
\begin{enumeratei}
\item\label{1.14Ki}
L'application $\Phi$ de \eqref{application-phi} est injective.
\item\label{1.14Kii}
L'application \[\Ker[H^i(k,M) \to H^i(L,M)] \to \Ker [H^i_{\on{nr}}(k(Y)/k,M) \to H^i_{\on{nr}}(L(Y)/L,M)]\] n'est pas
surjective.
\end{enumeratei}
Alors le groupe $A_{0}(X\times_{K}K(X))$ est non nul, en d'autres termes,
la $K$-variété $X$ n'est pas $\CH_{0}$-triviale.
\end{thm}

\begin{proof}
On suit la stratégie de la preuve de \cite[Th.\,1.12]{colliot2016hypersurfaces}. Soit $\eta$ le point générique de $Y$ et soit $B\coloneqq O_{\mc{X},\eta}$ son anneau local, qui est de dimension $1$. Comme la fibre $Y$ est géométriquement intègre, l'idéal maximal de $B$ est engendré par l'image $\nu$ d'une uniformisante de $A$. Comme $\pi$ est fidèlement plat, $\nu$ n'est pas diviseur de zéro, et donc $B$ est un anneau de valuation discrète de corps des fractions le corps des fonctions $K(X)$ de $X$ et de corps résiduel le corps des fonctions $k(Y)$ de~$Y$. On~considère le $B$-schéma ${\mathcal X}\times_{A}B$. Sa fibre générique est $X_{K(X)}$ et sa fibre spéciale est $Y_{k(Y)}$.

On a un homomorphisme de spécialisation
$\CH_{0}(X_{K(X)}) \to \CH_{0}(Y_{k(Y)})$
qui envoie le point générique de $X$ sur le point générique de $Y$ et le point $\sigma(K)_{K(X)}$ sur $b_{k(Y)}$; voir \cite{fulton1975rational} et \cite[\S 20.3]{fulton1998intersection}.

D'après la proposition \ref{prop1}, la classe du point générique de $Y$
dans $\CH_{0}(Y_{k(Y)})$ n'est pas
égale à celle de $b_{k(Y)} $,
et donc $A_{0}(X_{K(X)}) \neq 0$.
\end{proof}

On donne une condition suffisante pour l'injectivité de $\Phi$ qui suffira pour tous les exemples de la section \ref{paraex}. Une autre condition suffisante sera donnée dans le lemme~\ref{lemme-phi-isomorphisme-variante} et sera utilisée pour les exemples de la section \ref{dimsup}.

\begin{lemma}\label{lemme-phi-isomorphisme}
On se place dans la situation \ref{situation-phi}. Si le morphisme $S \to \Spec(k)$ se factorise par $\Spec(L) \to \Spec(k)$, alors l'application $\Phi$ est un isomorphisme.
\end{lemma}

\begin{proof}
On a un diagramme commutatif
\[
\begin{tikzcd}[row sep=2.7em, column sep=3.5em]
\CH_{0}(T_{E'}) \arrow[r] \arrow[d,->>,"\ts N_{E'/E}"] \arrow[ddd, bend left=66,"\ts p_*"] &
\CH_{0}(Z_{E'}) \arrow[r] \arrow[d,"\ts N_{E'/E}"] \arrow[ddd, bend left=66,"\ts p_*"] &
\CH_{0}(p^{-1}(W)_{E'}) \arrow[r] \arrow[d,"\ts N_{E'/E}"] \arrow[ddd, bend left=66,"\ts \wr"',"\ts p_*"] &
0 \\
\CH_{0}(T_{E}) \arrow[r] \arrow[d,"\ts p_*"] &
\CH_{0}(Z_{E}) \arrow[r] \arrow[d,"\ts p_*"] &
\CH_{0}(p^{-1}(W)_{E}) \arrow[r] \arrow[d,"\ts \wr"',"\ts p_*"] &
0 \\
\CH_{0}(S_{E}) \arrow[r] &
\CH_{0}(Y_{E}) \arrow[r] &
\CH_{0}(W_{E}) \arrow[r] &
0 \\
\CH_{0}(S_{E'}) \arrow[r] \arrow[u,->>,swap,"\ts N_{E'/E}"] &
\CH_{0}(Y_{E'}) \arrow[r] \arrow[u,swap,"\ts N_{E'/E}"] &
\CH_{0}(W_{E'}) \arrow[r] \arrow[u,swap,"\ts N_{E'/E}"] &
0
\end{tikzcd}
\]
où les suites exactes horizontales proviennent de la suite de localisation, et où on écrit $N_{E'/E}$ pour les applications normes.

Par hypothèse, le $k$-schéma $S$ admet une structure de $L$-schéma. Comme le corps~$k$ est de caractéristique zéro, l'inclusion diagonale $L\to L\otimes_kL$ est scindée, et donc la projection de $L$-schémas $S\times_kL=S\times_L(L\otimes_kL)\to S$ admet une section. Ainsi la flèche norme $N_{E'/E}\colon \CH_{0}(S_{E'}) \to \CH_{0}(S_{E})$ est surjective. Comme $T$ admet un morphisme vers $S$, il~admet aussi une structure de $L$-schéma, et le même argument montre que la flèche norme $N_{E'/E}\colon \CH_{0}(T_{E'}) \to \CH_{0}(T_{E})$ est surjective.
On déduit alors du diagramme ci-dessus que l'application $\Phi$ est un isomorphisme.
\end{proof}

\section{Variétés projectives et lisses non stablement rationnelles sur\nobreakspace le\nobreakspace corps\nobreakspace \texorpdfstring{$R$}{R} des séries de Puiseux réelles dont l'espace des \texorpdfstring{$R$}{R}-points est semi-algébriquement connexe}
\label{paraex}

\begin{thm}\label{generateur}
Soit $Y$ une $\R$-variété projective et géométriquement intègre. Soit $S\subset Y$ le lieu singulier de $Y$, avec structure réduite. Supposons:
\begin{enumeratea}
\item\label{generateura}
L'ensemble $Y(\R)$ des points réels est contenu dans le lieu lisse de $Y$.
\item\label{generateurb}
L'espace topologique $Y(\R)$ est connexe (et en particulier non vide).
\item\label{generateurc}
Le morphisme $S\to \on{Spec}(\R)$ se factorise par $\Spec(\C)\to\Spec(\R)$.
\item\label{generateurd}
L'application
\[\Br(\R) \to \Ker [\Br_{\on{nr}}(\R(Y)/\R) \to \Br_{\on{nr}}(\C(Y)/\C)]\] n'est pas
surjective.
\end{enumeratea}

Soit $\A^1_\R=\Spec( \R[t])$, soit $U\subset \A^1_\R$ un voisinage Zariski de $0$, soit $f\colon \mc{X}\to U$ un morphisme projectif fidèlement plat de $\R$-variétés, à fibre générique lisse, tel que $\mc{X}|_{t=0}\cong Y$. Alors, si $\eta\colon \Spec(\puiseux)\to \A^1_\R$ est le morphisme correspondant à l'inclusion $\R[t]\subset \puiseux$, la $\puiseux$-variété projective lisse $X\coloneqq \mc{X}\times_{U,\eta}\Spec(\puiseux)$ satisfait les propriétés suivantes.

\begin{enumeratei}
\item\label{generateuri}
L'espace semi-algébrique $X(\puiseux)$ est semi-algébriquement connexe.
\item\label{generateurii}
La $\puiseux$-variété $X$ n'est pas $\CH_0$-triviale, et en particulier
n'est pas stablement rationnelle.
\end{enumeratei}
\end{thm}

\skpt
\begin{proof}
\eqref{generateuri}
Ceci suit du théorème \ref{critere-connexe} et des propriétés \eqref{generateura} et \eqref{generateurb}.

\eqref{generateurii}
Si la $\puiseux$-variété $X$ est $\CH_0$-triviale,
alors il existe un entier $n>0$ tel que la variété
$\mc{X}\times_U \R(\!(t^{1/n})\!)$ sur le corps $\R(\!(t^{1/n})\!)$
soit $\CH_0$-triviale. Ceci résulte de la forme ``décomposition de la diagonale''
de la $\CH_0$-trivialité \cite[Prop.\,1.4]{colliot2016hypersurfaces}. Pour montrer \eqref{generateurii}, il~suffit donc de prouver que pour tout $n\geq 1$ la variété
$\mc{X}\times_U \R(\!(t^{1/n})\!)$ n'est pas $\CH_0$-triviale.

Soit $n\geq 1$ un entier, soit $A\coloneqq \R[\![t^{1/n}]\!]$ et soit $\mc{X}'\coloneqq \mc{X}\times_U A$, où le produit fibré est pris par rapport au morphisme évident $\Spec(A)\to U$. La~fibre spéciale du morphisme naturel $\pi\colon \mc{X}'\to \Spec(A)$ est isomorphe à $Y$. D'après \eqref{generateura} et \eqref{generateurb}, il~existe un $\R$-point lisse $y\in Y(\R)$, qui se relève en une section $\sigma$ de $\pi$ d'après le lemme de Hensel. Ceci, les propriétés \eqref{generateurc} et \eqref{generateurd} et le lemme \ref{lemme-phi-isomorphisme} nous garantissent que $\pi$ satisfait les hypothèses \eqref{1.14Ki} et \eqref{1.14Kii} (avec $i=2$ et $M=\mu_2$) du théorème \ref{1.14K}, ce qui entraîne que la fibre générique $(\mc{X}')_{\R(\!(t^{1/n})\!)}\cong\mc{X}_{\R(\!(t^{1/n})\!)}$ n'est pas $\CH_0$-triviale, comme voulu.
\end{proof}

Ce théorème va nous permettre de construire des exemples projectifs et lisses $X$ sur $\puiseux$ en partant d'exemples projectifs singuliers $Y$ sur $\R$. Les trois premiers exemples singuliers avaient été construits dans \cite{colliot2024certaines}.

\subsection{Fibrations en surfaces quadriques sur la droite projective dont toutes les fibres géométriques sont intègres}\label{contreexemple}

Soit $\P^1_\R=\on{Proj}(\R[u_0,u_1])$, vu comme recollement de $\A^1_\R$ avec coordonnée $u\coloneqq u_0/u_1$ et de $\A^1_\R$ avec coordonnée $v\coloneqq u_1/u_0$ au moyen du changement de variable $u=1/v$ sur $\P^1_\R\setminus\{0,\infty\}$. Soit $Y$ le recollement de la sous-variété de $\P^3_{\R} \times_{\R} \A^1_{\R}$ avec coordonnées $(a,b,c,d;u)$ donnée par l'équation
\[a^2+(1+u^2) b^2 - u (c^2+d^2)=0\]
avec la sous-variété de $\P^3_{\R} \times_{\R} \A^1_{\R}$ avec coordonnées $(a',b',c',d'; v)$ donnée par l'équation
\[a'^2+(1+v^2)b'^2-v(c'^2+d'^2)=0\]
au moyen du changement de variable
\[(a',b',c',d';v)=(a/u, b, c, d; 1/u).\] On a un morphisme surjectif $\pi\colon Y \to \P^1_\R$ donné par $(u_0:u_1)$.

\begin{lemma}\label{refX}
Soit $\pi \colon Y \to \P^1_{\R}$ la fibration en surfaces quadriques définie ci-dessus. Alors $Y$ satisfait les propriétés \eqref{generateura}--\eqref{generateurd} du théorème \ref{generateur}.
\end{lemma}
\begin{proof}
Les propriétés \eqref{generateura} et \eqref{generateurc} se vérifient en calculant explicitement le lieu singulier de $Y$, qui est formé d'un nombre fini de points dont aucun n'est réel. Pour montrer \eqref{generateurb} (la connexité), il~suffit de noter que l'application $Y(\R) \to \P^1(\R)$ a pour image $[0,\infty]$, qui est connexe, et que pour tout $s\in [0,\infty]$ la fibre en $s$ est connexe: pour $0<s<\infty$, les points réels des fibres sont des quadriques connexes, et pour $s=0,\infty$, ce sont des droites projectives. Pour \eqref{generateurd}, il~suffit de rappeler que, d'après \cite[Prop.\,13.1]{colliot2024certaines}, l'image de la classe de quaternions $(-1,u) \in \Br(\R(u))=\Br(\R(\P^1))$ dans $\Br(\R(Y))$ est non ramifiée, s'annule dans $\Br_{\on{nr}}(\C(Y)/\C)$, et n'est pas dans l'image de $\Br(\R)$.
\end{proof}

Soit $P(u)=a_0 + a_1u + a_2 u^2 \in \R[u]$, avec $a_0a_2\neq 0$.
On considère le sous-schéma fermé de $\P^3_{\R[t]} \times_{\R[t]} \A^1_{\R[t]}$ avec coordonnées $(a,b,c,d;u)$ donné par l'équation
\[a^2+(1+u^2)b^2 + (tP(u)- u)c^2 -ud^2=0.\]
On recolle ce schéma avec le sous-schéma fermé de $\P^3_{\R[t]} \times_{\R[t]} \A^1_{\R[t]}$ avec coordonnées $(a',b',c',d';v)$
donné par l'équation
\[a'^2+(1+v^2)b'^2 + (tQ(v)-v)c'^2 -vd'^2=0,\]
où $Q(v)\coloneqq v^2P(1/v)= a_0v^2 + a_1v + a_2$, via le changement de variables
\[(a',b',c',d';v)=(a/u, b, c, d; 1/u).\]
On obtient une fibration en quadriques $\mc{X} \to \P^1_{\R[t]}$. La~fibre spéciale en $t=0$ est la $\R$-variété $Y$. La~fibre générique sur le corps $\R(t)$
est une variété projective lisse munie d'une fibration en quadriques sur $\P^1_{\R(t)}$ dont toutes les fibres géométriques sont intègres et avec exactement $6$ fibres géométriques singulières.

Le théorème \ref{generateur} et le lemme \ref{refX} donnent maintenant :
\begin{thm}\label{prop-fibre-quadriques}
Soit $\mc{X}\to \P^1_{\R[t]}$ le morphisme défini ci-dessus.
La $\puiseux$-variété $X\coloneqq \mc{X}\times_{\R[t]} \puiseux$ est une variété projective lisse fibrée en surfaces quadriques sur $\P^1_{\puiseux}$ à fibres géométriques intègres. Elle n'est pas $\CH_0$-triviale et donc n'est pas stablement rationnelle. L'espace $X(\puiseux)$ est semi-algébriquement connexe.
\end{thm}

Soit $R$ un corps réel clos. Dans \cite{colliot2024certaines} et \cite{benoist2026rationality} on a étudié des fibrations en surfaces quadriques $X'\to \P^1_R$ d'équation affine
\begin{equation}\label{u-p(u)} x^2+y^2+z^2= u p(u)\end{equation}
avec $p(u)\in R[u]$ un polynôme séparable. Si $p(u)$ est strictement positif sur $R$, l'espace $X'(R)$ est semi-algébriquement connexe.

\begin{prop}
La fibre générique de la fibration $X \!\to\! \P^1_{\puiseux}$ du théorème~\ref{prop-fibre-quadriques} n'est isomorphe à la fibre générique d'aucune fibration $X'\to \P^1_{\puiseux}$ de modèle affine~\eqref{u-p(u)}.
\end{prop}
\begin{proof}
Soient plus généralement $k$ un corps de caractéristique différente de $2$ et $p\colon X \to \P^1_{k}$ une fibration en surfaces quadriques dont toutes les fibres
géométriques sont intègres. On~suppose qu'il existe au moins une fibre
géométrique singulière.
À une telle fibration on associe le revêtement discriminant $\Delta \to \P^1_{k}$.
La courbe~$\Delta$ est projective et lisse, et munie d'une classe
$\alpha \in \Br(\Delta)$ qui est obtenue de la façon suivante.
On prend une section lisse quelconque de la quadrique générique sur $k(\P^1)$
et on l'étend au corps $k(\Delta)$. On~obtient alors une quadrique sur
$k(\Delta)$ définie par une forme quadratique $\langle 1,-a,-b,ab\rangle$. La~classe $\alpha$ est la classe de l'algèbre de quaternions $(a,b)\in \on{Br}(k(\Delta))$. Pour plus de détails, on consultera
\cite[Th.\,2.5]{colliot1993groupe}.
L'hypothèse sur les fibres singulières
assure que l'on a $\alpha \in \Br(\Delta)$.
Pour les familles $X' \to \P^1_{k}$ considérées
dans \cite{colliot2024certaines} et \cite{benoist2026rationality}, la classe $\alpha$ est dans l'image
de $\Br(k) \to \Br(\Delta)$. De fait, pour une fibration d'équation affine
$x^2+y^2+z^2=P(u)$, avec $P(u)\in k[u]$, la classe $\alpha \in \Br (k(\Delta))$ est l'image de $(-1,-1) \in \Br(k)$.

Soit $R=\puiseux$. Dans le cas considéré au théorème
\ref{prop-fibre-quadriques}, avec $k=R$,
la courbe $\Delta$ est définie sur $R$ par l'équation
\[
 w^2= (1+u^2)(-tuP(u)+u^2).
\]
On prend la section de la quadrique générique donnée par $c=0$. La~classe
$\alpha \in \Br(\Delta)$ est donc l'image de $(u,-(1+u^2))=(-1,u) \in \Br(R(\P^1))$ dans
$\Br(R(\Delta))$. En un point $A \in \Delta(\R) \subset \Delta(R)$ avec $u>0$,
on a $\alpha(A)=0$. En un point $B \in \Delta(\R) \subset \Delta(R)$
avec $u<0$ on a $\alpha(B)\neq 0$. Ainsi $\alpha$ n'est pas dans
l'image de $\Br(R)$.
\end{proof}

\subsection{Intersections complètes lisses de deux quadriques dans \texorpdfstring{$\P^5$}{P5}}
Considérons l'espace projectif $\P^5_\R$ avec coordonnées homogènes $a,b,c,d,e,f$. La~variété du paragraphe \ref{contreexemple} contient l'ouvert affine défini par l'équation
\[a^2+1+u^2 - u (c^2+d^2)=0,\]
ou encore le système de deux équations affines
\begin{align*}
a^2+1+u^2 -uv &=0, \\
c^2+d^2-v&=0.
\end{align*}
En changeant les notations, on voit que
la variété ainsi définie est birationnelle
à l'intersection de
deux quadriques $Y\subset \P^5_{\R}$
donnée par les équations homogènes
\begin{align*}
a^2+ b^2+c^2 - cd&=0,\\
e^2+f^2-bd&= 0.
\end{align*}
Notons que $Y \subset \P^5_\R$ ne contient pas de droite réelle: en effet la section de $Y$ par $d=0$ est donnée par $a^2+ b^2+c^2=e^2+f^2=0$ et donc ne contient pas de point réel. (Cette remarque ne sera pas utilisée dans la suite.)

\begin{lemma}\label{speciale-int-quadriques}
Soit $Y\subset \P^5_\R$ l'intersection de deux quadriques définies ci-dessus. Alors~$Y$ satisfait les propriétés \eqref{generateura}--\eqref{generateurd} du théorème \ref{generateur}.
\end{lemma}

\begin{proof}
Le lieu singulier de $Y$ est formé des $4$ points non réels donnés par
\[
(a,b,c,d,e,f)= (\pm i, 0, 1,0,0,0),\, (0,0,0,0,\pm i, 1).
\]
Ceci entraîne \eqref{generateura} et \eqref{generateurc}. Comme $Y$ est birationnelle à la variété du paragraphe \ref{contreexemple}, la propriété \eqref{generateurb} suit des lemmes \ref{connexe-invariant} et \ref{refX}, et la propriété \eqref{generateurd} suit du lemme \ref{refX}.
\end{proof}

Soient $q_1, q_2$ deux formes quadratiques sur $\R$ en les 6 variables $(a,b,c,d,e,f)$. Supposons
que le système $q_1=q_2=0$ définisse une intersection complète lisse de deux quadriques dans $\P^5_{\R}$. On~sait qu'il suffit pour cela que le déterminant $\on{det}(q_1+tq_2)$ soit un polynôme séparable de degré $6$.

Soit $\mc{X}\subset \P^5_{\R[t]}$ défini par le système
\begin{align*}
a^2+ b^2+c^2 - cd +tq_1 &=0,\\
e^2+f^2 -bd +t q_2 &=0.
\end{align*}
Alors $\mc{X}|_{t=0}=Y$ et la fibre générique sur $\R(t)$ est une intersection complète lisse de deux quadriques. Le théorème \ref{generateur} et le lemme \ref{speciale-int-quadriques} donnent ici :

\begin{thm}\label{prop-intersection-quadriques}
La $\puiseux$-variété $X= \mc{X}\times_{\R[t]} \puiseux$ est une intersection complète lisse de deux quadriques dans $\P^5_{\puiseux}$ qui n'est pas $\CH_0$-triviale, et donc n'est pas
stablement rationnelle. L'espace $X(\puiseux)$ est semi-algébriquement connexe.\end{thm}

\begin{rmk}
On ne peut envisager une extension simple de l'argument ci-dessus aux intersections de deux quadriques en dimension supérieure. On~sait en effet
que sur tout corps $k$ de caractéristique zéro, pour toute variété $Y$ intersection complète géométriquement intègre non conique de deux quadriques dans $\P^n_k$ avec $n\geq 6$, et~tout modèle projectif et lisse $Z$ de $Y$, l'application
$\Br(k) \to \Br(Z)$ est surjective \cite[Th.\,3.8]{colliot1987intersectionsI}.
Par ailleurs, Hassett, Kollár et Tschinkel
\cite{hassett2022rationality} ont montré que toute intersection complète lisse de deux quadriques $X \subset \P^6_{\R} $
est $\R$-rationnelle dès que~$X(\R)$ est connexe.
\end{rmk}

\subsection{Hypersurfaces cubiques lisses dans \texorpdfstring{$\P^4$}{P4} }\label{exemplecubique}

La variété $Y$ du paragraphe \ref{contreexemple} est $\R$-birationnelle à l'hypersurface cubique singulière $Y \subset \P^4_{\R}$, géométriquement intègre, définie
par l'équation
\begin{equation}\label{Ycub}
b(a^2+b^2+c^2) - c(d^2+e^2)=0
\end{equation}
en coordonnées homogènes $a,b,c,d,e$.

\begin{lemma}\label{cubique}
L'hypersurface cubique singulière $Y \subset \P^4_{\R}$ définie ci-dessus satisfait les hypothèses \eqref{generateura}--\eqref{generateurd} du théorème \ref{generateur}.
\end{lemma}

\begin{proof}
Le lieu singulier de $Y$ est formé des $4$ points non réels donnés par
\[
(a,b,c,d,e) = (\pm i, 0,1,0,0),\, (0,0,0, \pm i, 1).
\]
Ceci entraîne \eqref{generateura} et \eqref{generateurc}. Comme $Y$ est birationnelle à la variété du paragraphe \ref{contreexemple}, la propriété \eqref{generateurb} suit des lemmes \ref{connexe-invariant} et \ref{refX}, et la propriété \eqref{generateurd} suit du lemme \ref{refX}.
\end{proof}

Soit $g(a,b,c,d,e)$ une forme cubique non singulière. Soit $\mc{X} \subset \P^4_{\R[t]}$ le schéma défini par l'équation
\[
b(a^2+b^2+c^2) - c(d^2+e^2) -tg(a,b,c,d,e)=0.
\]
On a $\mc{X}|_{t=0}=Y$. Le théorème \ref{generateur} et le lemme \ref{cubique} donnent ici :

\begin{thm}\label{prop-cubique}
La $\puiseux$-variété $X\coloneqq \mc{X}\times_{\R[t]} \puiseux$ est une hypersurface cubique lisse dans $\P^4_{\puiseux}$ qui n'est pas $\CH_0$-triviale, et donc n'est pas
stablement rationnelle. L'espace $X(\puiseux)$ est semi-algébriquement connexe.
\end{thm}

\begin{rmk}
Pour les hypersurfaces cubiques singulières sur $\R$, des résultats de non-rationalité ont été obtenus par Cheltsov, Tschinkel et Zhang \cite{cheltsov2024rationality}.
\end{rmk}

\subsection{Certaines fibrations en coniques sur \texorpdfstring{$\P^2$}{P2}.}

Soit $R$ un corps réel clos et soit $C$ sa clôture algébrique. Dans \cite{benoist2020clemens}, Benoist et Wittenberg considèrent les solides lisses $X$ fibrés en coniques sur
$\P^2_R$, donnés par des équations (voir ci-dessous des explications sur la notation)
\[
s^2+v^2=f(x,y,z)
\]
avec $f$ une forme homogène de degré $4$, positive,
définissant une courbe lisse dans~$\P^2_R$. Pour tout solide $X$ de cette forme, $X(R)$ est semi-algébriquement connexe (\cf remarque \ref{fibres-connexes}) et $X_C$ est $C$-rationnel. Benoist et Wittenberg montrent que $X$ n'est pas $R$-rationnel; voir \cite[Cor.\,3.6]{benoist2020clemens}. Nous allons montrer qu'il existe de tels solides lisses $X$ sur le corps $\puiseux$ qui ne sont pas $\CH_0$-triviaux, donc ne sont pas stablement rationnels, et pour lesquels $X(\puiseux)$ est semi-algébriquement connexe, et qui sont géométriquement rationnels.

Soit $A$ un anneau, soit $S\coloneqq \P^2_A$, soit $d\geq 0$ un entier et soit $F \in H^0(S, \mc{O}(2d))$ une section avec lieu d'annulation $C\subset S$. Soit $\P \coloneqq \P_S(\mc{O}(-d) \oplus \mc{O}(-d) \oplus \mathcal{O})$ et soit $p\colon \P\to S$ le morphisme de projection. D'après \cite[Ch.\,2, Prop.\,7.11(a)]{hartshorne} on a un isomorphisme \[p_* \mathcal{O}_{\P}(1) \simeq \mc{O}(-d) \oplus \mc{O}(-d) \oplus \mathcal{O}\] et on définit $u \in H^0(\P, \mathcal{O}_{\P}(1))$ comme l'image de la section $1\in A=H^0(S,\mc{O})$.
La~formule de projection nous donne un isomorphisme
\[ p_*(p^*\mc{O}(d) \otimes \mathcal{O}_{\P}(1)) \simeq \mathcal{O} \oplus \mathcal{O} \oplus \mc{O}(d),\]
et on définit $s, v \in H^0(\P, p^*\mc{O}(d) \otimes \mathcal{O}_{\P}(1))$ comme les images de $1\in H^0(S,\mc{O})$, vu~com\-me premier et deuxième facteur, respectivement. Les sections $s^2$, $v^2$ et $u^2F$ appartiennent alors à $H^0(\P,p^*\mc{O}(2d)\otimes \mc{O}_\P(2))$, et on pose
\[Q_{A,F}\coloneqq \{s^2+v^2=u^2F\}\subset \P.\]

\begin{lemma}\label{refBW}
Soient $q_1,q_2\in H^0(\P^2_\R, \mc{O}(2))$ deux formes quadratiques positives telles que les coniques lisses $q_1=0$ et $q_2=0$ dans $\P^2_\R$ soient transverses. Soit $Y\coloneqq Q_{\R, q_1q_2}$ et soit $\pi \colon Y \to \P^2_{\R}$ le morphisme de projection. Alors $Y$ satisfait les propriétés \eqref{generateura}--\eqref{generateurd} du théorème \ref{generateur}.
\end{lemma}

\begin{proof}
Le lieu singulier de $Y$ est donné par l'équation $q_{1}=q_{2}=s=v=0$ et donc consiste en un nombre fini de points fermés non réels. Ceci montre \eqref{generateura} et \eqref{generateurc}.

L'application $Y(\R) \to \P^2(\R)$ est surjective et ses fibres
sont connexes. Ainsi $Y(\R)$ est connexe. Ceci donne \eqref{generateurb}.

Montrons \eqref{generateurd}. La~$\C$-variété $Y_\C$ étant birationnelle à $s'v'= q_1q_2$, elle est rationnelle, donc en particulier $\Br_{\mathrm{nr}}(\C(Y)/\C)=0$. Soit $h\in H^0(\P^2_\R,\mc{O}(1))$ et, pour $i=1,2$, posons $p_i\coloneqq q_i/h^2\in \R(\P^2)$. Soit \[\alpha\coloneqq (-1,p_1)\in H^2(\R(\P^2), \Z/2).\] Par un calcul de résidus sur $q_1=0$ et $q_2=0$, on vérifie que cette classe ne provient pas d'une classe dans $\Br(\R)$, et qu'elle n'est pas égale à la classe $(-1, p_1p_2)$.
Puisque le noyau de l'application $H^2(\R(\P^2), \Z/2)\to H^2(\R(Y), \Z/2)$ est engendré par la classe $(-1, p_1p_2)$, on déduit que l'image $\beta$ de $\alpha$ dans $H^2(\R(Y), \Z/2)$ n'est pas nulle. Il suffit de montrer que $\beta$ est non ramifiée \cite[Prop.\,4.2.3(a)]{colliot1995birational}.

Soit $w$ une valuation discrète de $\R(Y)$, soit $\kappa(w)$ son corps résiduel et soit $c_w$ le centre de $w$ sur $\P^2_\R$. On~voit que si $c_w$
n'est pas sur la conique $q_1=0$, alors le résidu de~$\beta$ est nul. Si $c_w$ est le point générique de $q_1=0$ et $w(p_{1})$ est pair, alors le résidu de~$\beta$ est nul.
Si $c_w$ est le point générique de $q_1=0$ et $w(p_{1})$ est impair, alors on vérifie sur l'équation que $-1$ est un carré dans $\kappa(w)$. Si $c_w$ est un point fermé, alors $\kappa(w)=\C$, car la conique n'a pas de point réel, et $-1$ est encore un carré dans $\kappa(w)$.
\end{proof}

Soit $f=f(x,y,z)$ un polynôme homogène réel de degré $4$ définissant une courbe lisse dans $\P^2_\R$ et soit $F\coloneqq q_1q_2+tf\in H^0(\P^2_{\R[t]},\mc{O}(4))$. Soit $X$ le fibré en coniques sur~$\P^2_{\puiseux}$ donné par l'équation
\[s^2+v^2=F(x,y,z).\]
Plus précisément, on pose $X\coloneqq Q_{\puiseux,F}$. Considérons $\mc{X}\coloneqq Q_{\R[t],F}$ avec sa projection naturelle $\mc{X}\to \P^2_{\R[t]}$. On~a $\mc{X}|_{t=0}=Y$ et $\mc{X}\times_{\R[t]}\puiseux=X$. Un calcul facile montre que la $\puiseux$-variété $X$ est lisse. Soit $C=\C\{\!\{t\}\!\}$. Alors un ouvert affine de $X_C$ est donné par l'équation $s'v'=F(x,y,1)$, où $s'=s+iv$ et $v'=s-iv$, ce qui entraîne la $C$-rationalité de $X_C$.

Le théorème \ref{generateur} et le lemme \ref{refBW} donnent ici:

\begin{thm}\label{prop-fibres-coniques}
La $R$-variété projective, lisse et géométriquement rationnelle
$X\coloneqq \mc{X}\times_{\R[t]}R$, où $R$ est le corps des séries de Puiseux réelles en $t$, est un solide fibré en coniques sur $\P^2_R$ tel que $X(R)$ soit semi-algébriquement connexe mais la variété $X$ n'est pas $\CH_0$-triviale, et en particulier n'est pas stablement rationnelle.
\end{thm}

\section{La méthode de spécialisation revisitée, II}\label{methode-2}

\begin{lemma}\label{lemme-phi-isomorphisme-variante}
On se place dans la situation \ref{situation-phi} et on suppose de plus que les homomorphismes $p_*\colon \CH_0(T_{E})\to \CH_0(S_{E})$ et $p_*\colon \CH_0(T_{E'})\to \CH_0(S_{E'})$ sont des isomorphismes. Alors l'application $\Phi$ est un isomorphisme.
\end{lemma}

\begin{proof}
Considérons le diagramme commutatif
\begin{displaymath}
\begin{adjustbox}{max width=\textwidth}
\begin{tikzcd}[row sep=3em, column sep=3em]
\CH_1(p^{-1}(W)_{E'},1) \arrow[d,"\ts N_{E'/E}"] \arrow[ddd, bend left=66,"\ts \wr"',"\ts p_*"] \arrow[r] &
\CH_{0}(T_{E'}) \arrow[r] \arrow[d,"\ts N_{E'/E}"] \arrow[ddd, bend left=66,"\ts \wr"',"\ts p_*"] &
\CH_{0}(Z_{E'}) \arrow[r] \arrow[d,"\ts N_{E'/E}"] \arrow[ddd, bend left=66,"\ts p_*"] &
\CH_{0}(p^{-1}(W)_{E'}) \arrow[r] \arrow[d,"\ts N_{E'/E}"] \arrow[ddd, bend left=66,"\ts \wr"',"\ts p_*"] &
0 \\
\CH_1(p^{-1}(W)_E,1) \arrow[r] \arrow[d,"\ts \wr"',"\ts p_*"] &
\CH_{0}(T_{E}) \arrow[r] \arrow[d,"\ts \wr"',"\ts p_*"] &
\CH_{0}(Z_{E}) \arrow[r] \arrow[d,"\ts p_*"] &
\CH_{0}(p^{-1}(W)_{E}) \arrow[r] \arrow[d,"\ts \wr"',"\ts p_*"] &
0 \\
\CH_1(W_{E},1) \arrow[r] &
\CH_{0}(S_{E}) \arrow[r] &
\CH_{0}(Y_{E}) \arrow[r] &
\CH_{0}(W_{E}) \arrow[r] &
0 \\
\CH_1(W_{E'},1) \arrow[r] \arrow[u,swap,"\ts N_{E'/E}"] &
\CH_{0}(S_{E'}) \arrow[r] \arrow[u,swap,"\ts N_{E'/E}"] &
\CH_{0}(Y_{E'}) \arrow[r] \arrow[u,swap,"\ts N_{E'/E}"] &
\CH_{0}(W_{E'}) \arrow[r] \arrow[u,swap,"\ts N_{E'/E}"] &
0
\end{tikzcd}
\end{adjustbox}
\end{displaymath}
où les complexes horizontaux sont exacts et où les applications $p_*\colon \CH_0(T_E)\to \CH_0(S_E)$ et $p_*\colon \CH_0(T_{E'})\to \CH_0(S_{E'})$ sont des isomorphismes par hypothèse. Ici $\CH_1(-,1)$ désigne le groupe de Chow supérieur au sens de Bloch \cite{bloch1986algebraic}; voir \cite[Prop.\,1.3]{bloch1986algebraic} pour la covariance par morphismes propres et \cite[Th.\,3.1]{bloch1986algebraic} pour la suite de localisation. (On aurait pu utiliser le formalisme de Rost \cite{rost1996chow}: voir \cite[(3.4) \& p.\,365]{rost1996chow} pour la covariance par morphismes propres, et voir \cite[p.\,356]{rost1996chow} pour la suite de localisation.) On en déduit que dans le carré commutatif
\[
\begin{tikzcd}[column sep = large]
\CH_0(Z_{E'}) \arrow[d,"\ts \wr"',"\ts p_*"] \arrow[r,"\ts N_{E'/E}"] & \CH_0(Z_E) \arrow[d,"\ts \wr"',"\ts p_*"] \\
\CH_0(Y_{E'}) \arrow[r,"\ts N_{E'/E}"] & \CH_0(Y_E)
\end{tikzcd}
\]
les flèches verticales sont des isomorphismes, ce qui entraîne que l'application
\[\Phi\colon \CH_{0}(Z_{E})/ N_{E'/E}(\CH_{0}(Z_{E'})) \to \CH_{0}(Y_{E})/N_{E'/E}(\CH_{0}(Y_{E'}))\]
est un isomorphisme.
\end{proof}

\begin{thm}\label{generateur-variante}
Soit $Y$ une $\R$-variété projective et géométriquement intègre. Soit $S\subset Y$ le lieu singulier de $Y$, avec structure réduite. Soit $W\coloneqq Y\setminus S$ l'ouvert complémentaire de $S$, soit $p\colon Z\to Y$ une résolution des singularités telle que la restriction $p^{-1}(W)\to W$ soit un isomorphisme, et soit $T\coloneqq p^{-1}(S)$. Soit $i\geq 0$ un entier. Supposons:

\begin{enumeratea}
\item\label{generateur-variantea}
les applications
\[p_*\colon \CH_0(T_{\R(Y)})\to \CH_0(S_{\R(Y)}),\qquad p_*\colon \CH_0(T_{\C(Y)})\to \CH_0(S_{\C(Y)})\] sont des isomorphismes;
\item\label{generateur-varianteb}
l'espace topologique $Y(\R)$ est connexe;
\item\label{generateur-variantec}
l'application
\[H^i(\R,\Z/2) \to \Ker [H^i_{\on{nr}}(\R(Y)/\R,\Z/2) \to H^i_{\on{nr}}(\C(Y)/\C,\Z/2)]\] n'est pas surjective.
\end{enumeratea}

Soit $\A^1_\R=\Spec(\R[t])$, soit $U\subset \A^1_\R$ un voisinage Zariski de $0$, soit $f\colon \mc{X}\to U$ un morphisme projectif fidèlement plat de $\R$-variétés, à fibre générique lisse, tel que $\mc{X}|_{t=0}\cong Y$. Alors, si $\eta\colon \Spec(\puiseux)\to\A^1_\R$ est le morphisme correspondant à l'inclusion $\R[t]\subset \puiseux$, la $\puiseux$-variété projective lisse $X\coloneqq \mc{X}\times_{U,\eta}\Spec(\puiseux)$ satisfait les propriétés suivantes.

\begin{enumeratei}
\item\label{generateur-variantei}
L'espace semi-algébrique $X(\puiseux)$ est semi-algébriquement connexe.
\item\label{generateur-varianteii}
La $\puiseux$-variété $X$ n'est pas $\CH_0$-triviale, et en particulier
n'est pas stablement rationnelle.
\end{enumeratei}
\end{thm}

\begin{proof}
On suit la preuve du théorème \ref{generateur}, en remplaçant le lemme \ref{lemme-phi-isomorphisme} par le lemme \ref{lemme-phi-isomorphisme-variante}.
\end{proof}

\begin{rmk}\label{fibre-par-fibre}
On rappelle une condition \enquote{fibre par fibre} qui entraîne l'hypothèse du lemme \ref{lemme-phi-isomorphisme-variante} (et donc, pour $k=\R$, l'hypothèse \eqref{generateur-variantea} du théorème \ref{generateur-variante}) et qui sera utilisée dans la suite. Soient $V$ et $W$ deux variétés sur un corps $k$ et soit $f\colon V\to W$ un $k$-morphisme propre. On~suppose que, pour tout point $P$ du schéma $W$, de corps résiduel $k(P)$, la fibre $V_P$ est une $k(P)$-variété $\CH_0$-triviale. D'après \cite[Prop.\,1.8]{colliot2016hypersurfaces}, pour tout corps $F$ contenant $k$, l'application $f_*\colon \CH_0(V_F)\to \CH_0(W_F)$ est un isomorphisme.
\end{rmk}

\section{Exemples en dimension supérieure}\label{dimsup}

\subsection{Le modèle singulier} Soit $\P^1_\R=\on{Proj}(\R[u_0,u_1])$, vu comme recollement de~$\A^1_\R$ avec coordonnée $u\coloneqq u_0/u_1$ et de $\A^1_\R$ avec coordonnée $v\coloneqq u_1/u_0$ au moyen du changement de variable $u=1/v$ sur $\P^1_\R\setminus\{0,\infty\}$. Soit $Y\subset \P^1_\R\times \P^7_\R$ la variété définie par l'équation
\[X_0^2+X_1^2+X_2^2 + (u_0^2+u_1^2)X_3^2 - u_0u_1(X_4^2+X_5^2+X_6^2+X_7^2) =0,\]
où $X_0,\dots,X_7$ sont des coordonnées sur $\P^7_\R$. La~projection $Y\to \P^1_\R$ est une fibration en quadriques de dimension relative $6$. Un modèle affine de $Y$ est donné par l'équation
\begin{equation} \label{eq:affine}
x^2+y^2+z^2 + (1+u^2) - u(x'^2+y'^2+z'^2+t'^2) = 0.
\end{equation}

\begin{lemma}\label{deuxieme-modele}
L'espace topologique $Y(\R)$ est connexe et l'application
\[H^3(\R,\Z/2)\to \on{Ker}[H^3_{\on{nr}}(\R(Y),\Z/2)\to H^3_{\on{nr}}(\C(Y),\Z/2)]\] n'est pas surjective.
\end{lemma}

\begin{proof}
L'image de $Y(\R) \to \P^1(\R)$ est l'intervalle $u\geq 0$, et pour tout $u\geq 0$ la fibre $Y_u(\R)$ est connexe. Donc $Y(\R)$ est connexe.

La classe $(-1,-1,u)\in H^3(\R(Y),\Z/2)$ devient nulle dans $\C(Y)$. On~va montrer que cette classe est non ramifiée et qu'elle n'est pas dans l'image de $H^3(\R,\Z/2)$. Soit~$q$ la forme quadratique
$\langle 1,1,1,1+u^2, -u,-u,-u,-u\rangle$ sur le corps $\R(u)$.
La forme~$q$ est de rang $8$ et n'est pas semblable à une forme de Pfister, puisque son déterminant
n'est pas un carré. D'après Arason \cite[Satz 5.6]{arason1975cohomologische} (voir aussi \cite[Cor.\,9.6.2]{kahn2008formes}), l'application $H^3(\R(u),\Z/2) \to H^3(\R(u)(q),\Z/2)$ est injective.

Soit $\R\subset A\subset \R(Y)$ un anneau de valuation discrète
de corps des fractions $\R(Y)$ et de corps résiduel $\kappa$.
Soit $v$ la valuation associée à $A$. Supposons
$(-1,-1)_{\kappa} \neq 0$, c'est-à-dire que $-1$ n'est pas une somme de
deux carrés dans $\kappa$. Alors une somme non nulle de $4$ carrés dans
le complété de $\R(Y)$ en $v$ a sa valuation paire.
Donc $v(x'^2+y'^2+z'^2+t'^2)$ est paire.
Supposons $v(u)$ impaire. Dans le complété de $A$,
si~$v(u)>0$, alors $1+u^2$ est un carré. Si $v(u)<0$, alors
$1+u^2$ est aussi un carré dans le complété. Donc
$x^2+y^2+z^2 + 1+u^2$ est une somme de $4$ carrés dans le complété,
donc de valuation paire. Mais ceci contredit l'équation.
Ainsi, pour toute valuation $v$ telle que le corps résiduel $\kappa$
satisfasse $(-1,-1)_{\kappa} \neq 0$, la valuation $v(u)$ est paire.

On conclut que la classe $(-1,-1,u) \in H^3(\R(Y),\Z/2)$ a tous ses résidus nuls,
elle est donc non ramifiée. D'autre part, si $(-1,-1,u) \in H^3(\R(Y),\Z/2)$ provenait de $H^3(\R,\Z/2\Z)$ alors, l'application $H^3(\R(u),\Z/2)\to H^3(\R(Y),\Z/2)$ étant injective, la classe $(-1,-1,u) \in H^3(\R(u),\Z/2)$ proviendrait de $H^3(\R,\Z/2\Z)$. Ceci est impossible car la classe $(-1,-1,u) \in H^3(\R(u),\Z/2)$ est ramifiée.
\end{proof}

\subsection{Variétés lisses fibrées en quadriques}

Soit $\mc{E}=\oplus_{i=0}^7\mc{E}_i$ un fibré vectoriel sur $\P^1_\R$, où $\mc{E}_i\coloneqq \mc{O}(1)$ pour $0\leq i\leq 2$ et $\mc{E}_i=\mc{O}$ pour $3\leq i\leq 7$. Soit encore $\mc{L}\coloneqq \mc{O}(2)$. On~a $\mc{E}_i^\vee\otimes \mc{E}_i^\vee\otimes \mc{L}\cong \mc{O}$ si $0\leq i\leq 2$ et $\mc{E}_i^\vee\otimes \mc{E}_i^\vee\otimes \mc{L}\cong \mc{O}(2)$ si $3\leq i\leq 7$. On~dispose donc des sections globales
\begin{align*}
1&\in \Gamma(\P^1_\R,\mc{O})=\Gamma(\P^1_\R, \mc{E}_i^\vee\otimes \mc{E}_i^\vee\otimes \mc{L})\qquad &(0\leq i\leq 2), \\
u_0^2+u_1^2&\in \Gamma(\P^1_\R,\mc{O}(2))=\Gamma(\P^1_\R, \mc{E}_i^\vee\otimes \mc{E}_i^\vee\otimes \mc{L})\qquad &(i=3),\\
-u_0u_1& \in \Gamma(\P^1_\R,\mc{O}(2))=\Gamma(\P^1_\R, \mc{E}_i^\vee\otimes \mc{E}_i^\vee\otimes \mc{L})\qquad &(4\leq i\leq 7).
\end{align*}

Considérons la fibration en quadriques $Q \subset \P(\mc{E})$ définie par l'annulation de la section globale
\[
q \coloneqq Z_0^2 + Z_1^2 + Z_2^2 + (u_0^2+u_1^2)Z_3^2 - u_0u_1(Z_4^2 + Z_5^2 + Z_6^2 + Z_7^2)\in \Gamma(\P^1_\R,\mathrm{Sym}^2(\mc{E}^\vee) \otimes \mc{L}).
\]

Soit $U_1=\{u_1\neq 0\}\subset \P^1_\R$. Une base de $\mc{E}|_{U_1}$ en tant que $\mc{O}_{U_1}$-module libre est donnée par $f_0/u_1, f_1/u_1, f_2/u_1, f_3, f_4, f_5, f_6, f_7$, où $f_i$ est une section globale de $(\mc{E}_i)|_{U_1}$. Notons $X_i$ les coordonnées de fibre correspondantes. Alors $(Z_i)|_{U_1}=u_1X_i$ si $0\leq i\leq 2$ et $(Z_i)|_{U_1}=X_i$ si $3\leq i\leq 7$.
En substituant dans l'équation $q=0$, on obtient l'équation locale de $Q$ dans $U_1 \times \P^7_\R$:
\[
u_1^2(X_0^2+X_1^2+X_2^2) + (u_0^2+u_1^2)X_3^2 - u_0u_1(X_4^2+X_5^2+X_6^2+X_7^2) = 0.
\]
Ceci équivaut à
\begin{equation} \label{eq:local_S}
X_0^2+X_1^2+X_2^2 + (1+u^2)X_3^2 - u(X_4^2+X_5^2+X_6^2+X_7^2) = 0,
\end{equation}
où $u=u_0/u_1$. En particulier, la variété affine \eqref{eq:affine} est isomorphe à l'ouvert de Zariski de $Q$ où $u_1\neq 0$ et $Z_3\neq 0$, en identifiant
\[(x,y,z,x',y',z',t')=(X_0/X_3,X_1/X_3,X_2/X_3,X_4/X_3,X_5/X_3,X_6/X_3,X_7/X_3).\]

Soit $U_0=\{u_0\neq 0\}\subset \P^1_\R$. Un argument analogue nous donne l'équation locale de~$Q$ dans $U_0 \times \P^7_\R$:
\[
u_0^2(Y_0^2+Y_1^2+Y_2^2) + (u_0^2+u_1^2)Y_3^2 - u_0u_1(Y_4^2+Y_5^2+Y_6^2+Y_7^2) = 0,
\]
où $(Z_i)|_{U_0}=u_0Y_i$ si $0\leq i\leq 2$ et $(Z_i)|_{U_0}=Y_i$ si $3\leq i\leq 7$.
Ceci équivaut à
\begin{equation} \label{eq:local_T}
Y_0^2+Y_1^2+Y_2^2 + (1+v^2)Y_3^2 - v(Y_4^2+Y_5^2+Y_6^2+Y_7^2) = 0,
\end{equation}
où $v=u_1/u_0$.

\begin{lemma}\label{fibration-quadriques-a-c}
Il existe une résolution des singularités $\tilde{Q}\to Q$ qui est un isomorphisme en dehors du lieu singulier de $Q$ telle que les hypothèses \eqref{generateur-variantea}, \eqref{generateur-varianteb} et \eqref{generateur-variantec} (avec $i=3$) du théorème \ref{generateur-variante} sont satisfaites pour $Y=Q$, $Z=\tilde{Q}$ et la résolution $\tilde{Q}\to Q$.
\end{lemma}

\begin{proof}
La $\R$-variété $Q$ est birationnelle à la $\R$-variété du lemme \ref{deuxieme-modele}. La~propriété \eqref{generateur-variantec} en résulte.

Le lieu singulier $S_1$ de $Q|_{U_1}$ est donné par les équations
\begin{align*}
X_0=X_1=X_2=(1+u^2)X_3=uX_4=uX_5=uX_6=uX_7&=0, \\
2uX_3^2-(X_4^2+X_5^2+X_6^2+X_7^2)&=0,
\end{align*}
et donc coïncide avec
\[S_1=\{u=X_0=X_1=X_2=X_3=X_4^2+X_5^2+X_6^2+X_7^2=0\}.\]
Le lieu singulier $S_0$ de $Q|_{U_0}$ est donné par
\[S_0=\{v=Y_0=Y_1=Y_2=Y_3=Y_4^2+Y_5^2+Y_6^2+Y_7^2=0\}.\]
Le lieu singulier $S$ de $Q$ est donc l'union disjointe de $S_0$ et $S_1$. En particulier $Q(\R)$ est contenu dans le lieu non singulier de $Q$. En particulier, on obtient \eqref{generateur-varianteb}: la~$\R$\nobreakdash-variété~$Q$ est birationnelle à la variété du lemme \ref{deuxieme-modele} et donc le lemme \ref{connexe-invariant} entraîne que $Q(\R)$ est connexe.

Les équations \eqref{eq:local_S} et \eqref{eq:local_T} montrent que le fibré en quadriques $Q\to \P^1_\R$ est admissible au sens de \cite[Déf.\,3.1]{colliot1993groupe}, c'est-à-dire, sa fibre générique est lisse et pour tout point fermé $P\in \P^1_\R$, la restriction de $Q\to \P^1_\R$ à $\Spec(O_P)$ est le lieu des zéros d'une équation homogène $\sum_{i=1}^8 a_iT_i^2$ où $0\leq v_P(a_i)\leq 1$ pour tout $i$ et $v_P(a_i)=0$ pour $1\leq i\leq 4$. D'après \cite[Th.\,3.3(a)]{colliot1993groupe}, une résolution des singularités $\tilde{Q}\to Q$ est obtenue en éclatant la variété $Q$ le long de son lieu singulier. De plus, pour tout point singulier $P\in Q$ la fibre du morphisme $\tilde{Q}\to Q$ en $P$ est donnée par une hypersurface quadrique lisse avec un $k(P)$-point rationnel (\cf \cite[bas de p.\,488]{colliot1993groupe}). Compte tenu de la remarque \ref{fibre-par-fibre}, ceci entraîne \eqref{generateur-variantea}.
\end{proof}

Soient $q_3(u),\dots,q_7(u)\in \R[u]$ des polynômes de degré $2$, sans racines doubles et tels que $\on{pgcd}(q_i,q_j)=1$ pour tout $3\leq i<j\leq 7$. Posons
\[p_i(u_0,u_1)\coloneqq u_1^2q_i(u_0/u_1)\in \Gamma(\P^1_\R,\mc{O}(2))=\Gamma(\P^1_\R,\mc{E}_i^\vee\otimes \mc{E}_i^\vee\otimes \mc{L})\qquad (3\leq i\leq 7).\]
Soit
\[q'\coloneqq Z_0^2+Z_1^2+Z_2^2+p_3Z_3^2+p_4Z_4^2+p_5Z_5^2+p_6Z_6^2+p_7Z_7^2\in \Gamma(\P^1_\R, \on{Sym}^2(\mc{E}^\vee)\otimes L),\]
et soit $Q'\coloneqq \{q'=0\}\subset \P(\mc{E})$. L'ouvert $(Q')_{U_1}\subset U_1\times \P^7_\R$ est donné par l'équation
\[X_0^2+X_1^2+X_2^2+q_3(u)X_3^2+q_4(u)X_4^2+q_5(u)X_5^2+q_6(u)X_6^2+q_7(u)X_7^2=0,\]
où $(Z_i)|_{U_1}=u_1X_i$ si $0\leq i\leq 2$ et $(Z_i)|_{U_1}=X_i$ si $3\leq i\leq 7$.
Le lieu singulier de $(Q')_{U_1}$ est donc donné par
\begin{align*}
X_0=X_1=X_2=q_3(u)X_3=q_4(u)X_4=q_5(u)X_5=q_6(u)X_6=q_7(u)X_7&=0,\\
(\partial q_3)(u)X_3^2+(\partial q_4)(u)X_4^2+(\partial q_5)(u)X_5^2+(\partial q_6)(u)X_6^2+(\partial q_7)(u)X_7^2&=0,
\end{align*}
où $\partial q_i$ est la dérivée de $q_i$ en $u$. Soit $P\in (Q')_{U_1}$ un point singulier. D'après les hypothèses faites sur les $q_i$, il~existe au plus un $j$ tel que $q_j(u(P))=0$, ce qui entraîne $X_i=0$ en $P$ pour tout $i\neq j$. Comme $q_j$ n'a pas de racine double, on a $(\partial q_j)(u(P))\neq 0$, et la dernière équation donne alors $X_j=0$ en $P$, contradiction. Comme les $q_i$ sont de degré $2$, la fibre de $Q'$ en $(1:0)$ est lisse. On~conclut que la $\R$-variété $Q'$ est lisse.

On considère la famille $\mc{X}\to \on{Spec}(\R[t])$ définie par l'équation $q+tq'=0$ dans $\P(\mc{E})\times_\R\R[t]$. La~fibre du morphisme $\mc{X}\to \on{Spec}(\R[t])$ en $t=0$ est $Q$ et la fibre générique est lisse. On~pose $X\coloneqq \mc{X}\times_{\R[t]}\puiseux$. Le théorème \ref{generateur-variante} et le lemme \ref{fibration-quadriques-a-c} nous donnent:

\begin{thm}\label{fibres-quadriques-6}
La $\puiseux$-variété $X$ est une variété lisse, projective, de dimension~$7$, fibrée en quadriques sur $\P^1_{\puiseux}$, non $\CH_0$-triviale, et telle que $X(\puiseux)$ soit semi-algébriquement connexe.
\end{thm}

\begin{rmk}\label{autrescontrex} Au théorème \ref{prop-fibre-quadriques} nous avons construit une variété
lisse, projective, de dimension $3$, fibrée en quadriques sur $\P^1_{\puiseux}$ avec les mêmes propriétés.
On part là d'une fibration
en quadriques $Y \to \P^1_\R$ dont la fibre générique est définie par une forme quadratique $q$ de rang $r=4$ sur $F=\R(\P^1)$
dont le déterminant n'est pas un carré. On~suppose en outre
que $Y$ possède un $\R$-point lisse.
On utilise
l'injectivité de $H^2(F,\Z/2) \to H^2(F(q), \Z/2)$,
qui résulte de l'hypothèse sur le déterminant,
et une classe non constante dans $H^2_{\on{nr}}(F(q)/\R, \Z/2)$.

Pour toute forme quadratique $q$ non dégénérée de rang $r\geq 5$
l'application \[
H^2(F,\Z/2) \to H^2_{\on{nr}}(F(q)/F, \Z/2)
\]
est un isomorphisme.
Pour $r\geq 5$, ceci implique que l'application
\[
H^2(\R,\Z/2) \to H^2_{\on{nr}}(F(q)/\R, \Z/2)
\]
est un isomorphisme.
On ne peut donc utiliser le groupe
$H^2(-, \Z/2)$ pour faire une construction analogue
avec $X$ de dimension plus grande que 3.

La construction du théorème \ref{fibres-quadriques-6}
part d'une fibration
en quadriques $Y \to \P^1_\R$ dont la fibre générique est définie par une forme quadratique de rang $r= 8$ sur $F=\R(\P^1)$
dont le déterminant n'est pas un carré, donc qui n'est pas semblable à une 3-forme de Pfister.
On utilise le fait que pour une telle forme~$q$ l'application
$H^3(F,\Z/2) \to H^3(F(q), \Z/2)$ est injective.

Plus généralement, d'après Arason \cite[Satz 5.6]{arason1975cohomologische},
cette application est injective si la forme $q$
est de rang $r\geq 5$ et n'est pas une 3-voisine de Pfister, \ie
n'est pas semblable à une sous-forme d'une 3-forme
de Pfister $\langle \!\langle a,b,c \rangle \!\rangle$.
Par ailleurs, on sait \cite[Th.\,10.2.4]{kahn2008formes} que pour $q$ avec $r\geq 3$, l'application
\[
H^3(F, \Q/\Z(2)) \to H^3_{\on{nr}}(F(q)/F, \Q/\Z(2))
\]
est surjective sauf si $r=6$ et $q$
est une forme d'Albert anisotrope.
On déduit de ces résultats et d'un calcul de résidus
que l'application $H^3(\R, \Z/2) \to H^3_{\on{nr}}(F(q)/\R, \Z/2)$
est surjective si $r\geq 9$.

On voit ainsi qu'outre les exemples que nous avons développés avec $r=4$ en utilisant $H^2$ et $r=8$ en utilisant $H^3$, on pourrait essayer d'utiliser $H^3(-, \Z/2)$ pour construire des exemples avec $r=5,6,7$ et $q$ non voisine d'une 3-forme de Pfister.
\end{rmk}

\subsection{Intersections de deux quadriques lisses}

Soit $\A^9_\R$ défini par les coordonnées $(x, y, z, u, v, x', y', z', t')$ et soit $U$ la variété réelle affine définie par les équations:
\begin{align}
x^2+y^2+z^2+1+u^2 - uv &= 0, \label{eq:U1} \\
x'^2+y'^2+z'^2+t'^2 - v &= 0. \label{eq:U2}
\end{align}
Soient $[X_0: X_1: X_2: X_3: X_4: X_5: X_6: X_7: X_8: X_9]$ des coordonnées homogènes sur $\P^9_\R$. On~identifie $\A^9_\R$ à $U_0=\{X_0\neq 0\}\subset \P^9_\R$ en posant
\[(x,y,z,u,v,x',y',z',t')=(X_1/X_0,X_2/X_0,\dots, X_9/X_0).\]
Définissons $Y$ par:
\begin{align}
X_0^2+X_1^2+X_2^2+X_3^2+X_4^2 - X_4X_5 &= 0, \label{eq:F1} \\
X_6^2+X_7^2+X_8^2+X_9^2 -X_0X_5 &= 0. \label{eq:F2}
\end{align}
Donc $U=Y\cap U_0$.

\begin{lemma}\label{int-2-quadriques-a-c}
Pour la $\R$-variété $Y$ ci-dessus, il~existe une résolution des singularités $p\colon Z\to Y$ qui est un isomorphisme en dehors du lieu singulier de $Y$ telle que les hypothèses \eqref{generateur-variantea}, \eqref{generateur-varianteb} et \eqref{generateur-variantec} (avec $i=3$) du théorème \ref{generateur-variante} sont satisfaites.
\end{lemma}

\begin{proof}
La $\R$-variété $Y$ est birationnelle à la $\R$-variété du lemme \ref{deuxieme-modele}. La~propriété \eqref{generateur-variantec} en résulte.

La matrice jacobienne associée au système d'équations \eqref{eq:F1}-\eqref{eq:F2} est :
\[
 \begin{pmatrix} 2X_0 & 2X_1 & 2X_2 & 2X_3 & 2X_4-X_5 & -X_4 & 0 & 0 & 0 & 0 \\ X_5 & 0 & 0 & 0 & 0 & X_0 & -2X_6 & -2X_7 & -2X_8 & -2X_9 \end{pmatrix}. 
\]
Soit $P$ un point singulier de $Y$. Si $X_5\neq 0$ en $P$, alors
\[X_1=X_2=X_3=X_4=2X_4-X_5=2X_0^2+X_4X_5=0,\] ce qui contredit \eqref{eq:F1}. Donc $X_5=0$ en $P$. Si $X_0\neq 0$ en $P$, alors \[0=X_4X_5=-2X_0^2\neq 0,\] contradiction. Donc $X_0=X_5=0$ en $P$. La~forme de la matrice jacobienne entraîne alors que l'on a soit $X_1=X_2=X_3=X_4=0$ en $P$ (avec $X_6,X_7,X_8,X_9$ arbitraires) soit $X_6=X_7=X_8=X_9=0$ (avec $X_1,X_2,X_3,X_4$ arbitraires). On~déduit que le lieu singulier $S\subset Y$, avec la structure réduite, est l'union disjointe de
\begin{align*}
S_1\coloneqq \{ X_0=X_1=X_2=X_3=X_4=X_5=0, \quad X_6^2+X_7^2+X_8^2+X_9^2=0\},\\
S_2\coloneqq\{ X_0=X_5=X_6=X_7=X_8=X_9=0, \quad X_1^2+X_2^2+X_3^2+X_4^2=0\}.
\end{align*}
En particulier, $Y(\R)$ est contenu dans le lieu lisse de $Y$ et connexe, ce qui, par le lemme \ref{connexe-invariant}, donne \eqref{generateur-varianteb}.

Soit $p\colon Z\to Y$ l'éclatement de $Y$ en $S$: on veut montrer que $p$ est une résolution des singularités satisfaisant la propriété \eqref{generateur-variantea}. On~montrera le résultat plus précis suivant: la $\R$-variété $Z$ est lisse et pour tout $P\in S$ la fibre de $p$ en $P$ est une quadrique projective lisse avec un $k(P)$-point (d'après la remarque \ref{fibre-par-fibre}, ceci entraîne \eqref{generateur-variantea}). Pour tout $0\leq i\leq 9$, soit $U_i\coloneqq \{X_i\neq 0\}\subset \P^9_\R$. Pour tout point $P\in S_1$ il existe $i$ avec $6\leq i\leq 9$ tel que $P\in U_i$. Par symétrie, on peut supposer $i=9$. L'ouvert $Y\cap U_9\subset \A^9_\R$ est donné par les équations
\begin{align*}
x_0^2+x_1^2+x_2^2+x_3^2+x_4^2-x_4x_5 &= 0,\\
x_6^2+x_7^2+x_8^2+1-x_0x_5&= 0,
\end{align*}
où $x_i\coloneqq X_i/X_9$ sont les coordonnées de $U_9=\A^9_\R$.

Comme on a $P\in S_1$, on a $x_0=0$ en $P$ et donc la deuxième équation entraîne l'existence de $j$ avec $6\leq j\leq 8$ tel que $x_j\neq 0$ en $P$. Par symétrie, on peut supposer $j=8$. Soit $\A^9_\R$ un autre espace affine, avec coordonnées $y_0,\dots,y_8$. On~a un morphisme
\[f\colon \A^9_\R\to \A^9_\R,\qquad (x_0,\dots,x_8)\mapsto (x_0,\dots,x_7,x_8^2).\]
Comme $x_8\neq 0$ en $P$, la restriction de $f$ à $Y\cap U_9\subset \A^9_\R$ est étale au voisinage de $P$. De plus, si $W\subset \A^9_\R$ est donné par
\begin{align*}
y_0^2+y_1^2+y_2^2+y_3^2+y_4^2-y_4y_5&=0,\\
y_6^2+y_7^2+y_8+1-y_0y_5&=0,
\end{align*}
alors $f^{-1}(W)=Y\cap U_9$. On~note que $W$ est isomorphe à la sous-variété de $\A^8_\R$, avec coordonnées $z_j$ pour $0\leq j\leq 7$, donnée par l'équation
\[z_0^2+z_1^2+z_2^2+z_3^2+z_4^2-z_4z_5=0,\]
ce qui est isomorphe à $\tilde{C}\times \A^2_\R$, où $\tilde{C}$ est le cône affine sur une quadrique projective lisse~$C$ avec un $\R$-point. En particulier, l'éclaté de la variété $\tilde{C}\times \A^2_\R$ en son lieu singulier $\{0\}\times \A^2_\R$ est une $\R$-variété lisse, et le diviseur exceptionnel est donné par la deuxième projection $C\times \A^2_\R\to \A^2_\R$. On~déduit que tout point de $p^{-1}(S_1)$ est un point lisse de~$Z$ et que pour tout $P\in S_1$ la fibre de $p$ en $P$ est une quadrique projective lisse avec un $k(P)$-point. On~peut traiter $S_2$ de façon similaire. D'après la remarque~\ref{fibre-par-fibre}, ceci établit
la propriété \eqref{generateur-variantea}.
\end{proof}

Soient $q_1$ et $q_2$ deux formes quadratiques sur $\R$ en les $10$ variables $X_0,\dots,X_9$, telles que le système $q_1=q_2=0$ définisse une intersection complète lisse de deux quadriques dans $\P^9_{\R}$.

Soit $\mc{X}\subset \P^9_{\R[t]}$ défini par le système
\begin{align*}
X_0^2+X_1^2+X_2^2+X_3^2+X_4^2 - X_4X_5 +tq_1=0, \\
X_6^2+X_7^2+X_8^2+X_9^2 -X_0X_5 + tq_2=0.
\end{align*}
Alors $\mc{X}|_{t=0}=Y$ et la fibre générique sur $\R(t)$ est une intersection complète lisse de deux quadriques dans $\P^9_{\R(t)}$.
On pose $X\coloneqq \mc{X}\times_{\R[t]}\puiseux$.
Le théorème \ref{generateur-variante} et le lemme~\ref{int-2-quadriques-a-c} donnent ici :

\begin{thm}\label{intersection-quadriques-p9}
L'intersection lisse de deux quadriques $X$ dans $\P^9_{\puiseux}$ n'est pas $\CH_0$-triviale, et l'espace $X(\puiseux)$ est semi-algébriquement connexe.
\end{thm}

\begin{rmk}
Au théorème \ref{prop-intersection-quadriques} nous avons construit des intersections complètes lisses $X$ de deux
quadriques dans $\P^5_{\puiseux}$ pour lesquelles $X(\puiseux)$ est semi-algébriquement connexe et $X$ n'est pas $\CH_0$-triviale. On~peut se demander s'il existe de tels exemples dans $\P^n_{\puiseux}$ pour d'autres valeurs $n\geq 6$, $n\neq 9$.
On pourrait essayer d'en construire pour
$n=6,7,8$ en partant de la remarque \ref{autrescontrex}.
Il convient cependant de se rappeler que d'après \cite{hassett2022rationality}
de tels exemples n'existent pas sur $\R$ pour $n=6$.
\end{rmk}

\section{Calcul d'invariants cohomologiques non ramifiés pour les exemples}\label{calcul}

\begin{prop}\label{br-surj}
Soit $k$ un corps de caractéristique zéro.
Soit $X$ une $k$-variété projective lisse
géométriquement intègre. Supposons que $X$ est $k$-birationnelle à une $k$-variété
de l'un des types suivants :

\begin{enumeratea}
\item\label{br-surja}
intersection complète lisse de deux quadriques dans $\P^n_{k}$ avec $n \geq 5$;
\item\label{br-surjb}
hypersurface cubique lisse dans $\P^n_{k}$ avec $n \geq 4$;
\item\label{br-surjc}
modèle projectif et lisse d'une $k$-variété géométriquement intègre $Y$ munie d'une
fibration en quadriques $Y \to \P^1_k$ à fibre générique lisse et dont toutes les fibres ont une composante géométriquement intègre de multiplicité $1$.
\end{enumeratea}

Alors $\Br(k) \to \Br(X)$ est surjectif.
\end{prop}

\begin{proof}

Dans les cas \eqref{br-surja} et \eqref{br-surjb}, cette flèche est un isomorphisme \cite[Th.\,8.3.2]{colliot2021brauer}.
Dans le cas \eqref{br-surjc}, on commence par observer que la flèche
\[\Br(k(\P^1)) \to \Br_{\on{nr}}(k(Y)/k(\P^1))= \Br(Y_{\eta})\] est surjective car la fibre générique $Y_{\eta}/k(\P^1)$
est une quadrique lisse. Une classe $\alpha$ dans $\Br(X) \subset \Br(k(X))= \Br(k(Y))$
est donc image d'une classe $\beta \in \Br(k(\P^1))$. Pour tout point fermé $P\in \P^1_k$, la fibre $Y_P$ admet une composante irréductible $Z_P\subset Y_P$ de multiplicité $1$. La~comparaison du résidu de $\beta$ en $P$ et du résidu de $\alpha$ au point générique de $Z_P$ (voir \cite[Cor.\,11.1.6]{colliot2021brauer}) donne que le résidu de $\beta$ en $P$ est trivial. Comme ceci vaut pour tout point fermé $P$ de $\P^1_k$, on conclut que $\alpha$ est image d'un élément $\gamma \in \Br(\P^1_{k})=\Br(k)$.
\end{proof}

\Subsection{Un théorème de Witt}

\begin{thm}\label{wittgen}
Soit $R$ un corps réel clos, et soit $K=R(\Gamma)$ le corps des fonctions d'une $R$-courbe
$\Gamma$ géométriquement intègre. Soient $q_{1}$ et $q_{2}$ deux formes quadratiques non dégénérées
sur le corps $K$. Supposons $1 \leq {\rm rang}(q_{1}) < {\rm rang}(q_{2})$ et $3 \leq {\rm rang}(q_{2})$.
Si après évaluation en presque tout point $c \in \Gamma(R)$, la forme quadratique
$q_{1,c}$ est isomorphe à une sous-forme de $q_{2,c}$, alors la $K$-forme $q_{1}$
est isomorphe à une sous-forme de $q_{2}$. En particulier, si une forme quadratique $q$
non dégénérée sur $K$ de rang au moins $3$ se spécialise en une forme $q_{c}$
isotrope sur $R$ pour presque tout $c$, alors cette forme sur $K$ a un zéro non trivial.
\end{thm}

\begin{proof}
Dans le cas $R=\R$, ceci se déduit en combinant deux résultats de Witt :
\cite[Satz 22]{witt1937theorie} sur l'isotropie des formes quadratiques sur $R(\Gamma)$, ce qui correspond au cas du plan hyperbolique $q_{1}(x,y)=x^2-y^2$, et \cite[Satz 4]{witt1937theorie} sur la simplification des formes quadratiques.

Soit maintenant $R$ un corps réel clos quelconque. Soit $\Omega$ le spectre réel de $K$, c'est-à-dire l'espace topologique des ordres de $K$ (\cf \cite[(0.4.1)]{scheiderer1994real}) et, pour tout $\xi\in \Omega$, soit $K_\xi$ la clôture réelle de $K$ correspondante à $\xi$. L'hypothèse entraîne l'existence d'un sous-ensemble dense $S$ de $\Omega$ tel que, pour tout $\xi\in S$, la forme quadratique $(q_1)_{K_\xi}$ est isomorphe à une sous-forme de $(q_2)_{K_\xi}$; voir \cite[p.\,308]{scheiderer1996hasse}. Le $K$-schéma des représentations
de $q_{1}$ par $q_{2}$, qui est un espace homogène $P$ sous le groupe linéaire
connexe $\mathrm{SO}(q_{2})$, a donc un $K_\xi$-point pour tout $\xi\in S$. D'après \cite[Cor.\,6.2]{scheiderer1996hasse}, on déduit $P(K)\neq \emptyset$. Donc $q_1$ est isomorphe à une sous-forme de $q_2$ sur $K$.
\end{proof}

\begin{rmk}
Cet énoncé est en défaut si le rang de $q_{1}$ est égal au rang de $q_{2}$.
\end{rmk}

\subsection{Fibrations en quadriques et intersections de deux quadriques}

On calcule $H^3_{\on{nr}}(-,\Z/2)$ pour certaines fibrations en quadriques. En utilisant le théorème suivant, on en déduit des résultats analogues pour certaines intersections lisses de deux quadriques avec un point rationnel.

\begin{thm}\label{theorem3.2}
Soit $k$ un corps parfait de caractéristique différente de $2$, soit \hbox{$X\subset \P^n_k$} une intersection complète lisse de deux quadriques telle que $X(k)$ soit non vide. À tout $k$-point $P$ de $X$, on associe une fibration en quadriques $X'\to\P^1_k$, où $X'$ est lisse, telle que toute fibre géométrique ait au plus un point singulier (et donc soit intègre si $n\geq 5$) et telle que $X'$ soit $k$-birationnelle à $X$.
\end{thm}

\begin{proof}
Voir \cite[Th.\,3.2, (3.4) p.\,61, Rem.\,1.13.1]{colliot1987intersectionsI}.
\end{proof}

Nous utiliserons aussi le théorème suivant.
\begin{thm}\label{KRS}
Soit $k$ un corps de caractéristique différente de $2$.
Soit $q$ une forme quadratique non dégénérée sur $k$, de rang $r\geq 3$. Soit $k(q)$ le corps des fonctions de la quadrique projective définie par $q=0$.
Soit $i \geq 0$ un entier. Dans chacun des cas suivants, l'application
\[
H^{i}(k,\Q/\Z(i-1)) \to H^{i}_{\on{nr}}(k(q), \Q/\Z(i-1))
\]
est surjective:
\begin{enumeratei}
\item\label{KRSi}
$i=2$ et $r\geq 3$;

\item\label{KRSii}
$i=3$ et $r\neq 6$;

\item\label{KRSiii}
$i=3$, $r=6$ et $q$ n'est pas une forme d'Albert;

\item\label{KRSiv}
$i=4$ et $r\leq 5$.
\end{enumeratei}
\end{thm}

\begin{proof} Le cas $i=2$ est classique. Les énoncés pour $i=3$ sont dus à Kahn, Rost et Sujatha \cite[Th.\,5, Cor.\,10(2)]{kahn1998unramified}. L'énoncé \eqref{KRSiv} est dû à Kahn et Sujatha \cite[Th.\,3]{kahnsuja1998unramified}.
\end{proof}

\begin{prop}\label{H3nrRpincquad}
Soit $R$ un corps réel clos.
Soit $X$ une $R$-variété projective et lisse munie d'un
morphisme $X \to \P^1_{R}$ dont la fibre générique est une quadrique de dimension $d$. Supposons $X(R)$ semi-algébriquement connexe.
\begin{enumeratea}
\item\label{H3nrRpincquada}
Si $d=2$, l'application $H^3(R,\Z/2) \to H^3_{\on{nr}}(R(X)/R,\Z/2)$ est un isomorphisme.

\item\label{H3nrRpincquadb}
Si $d=3$, l'application
$H^4(R,\Z/2) \to H^4_{\on{nr}}(R(X)/R,\Z/2)$
est un isomorphisme.
\end{enumeratea}
\end{prop}

\begin{proof}
Démontrons d'abord l'énoncé \eqref{H3nrRpincquada}.
Notons $C=R(\sqrt{-1})$. La~fibration $X_{C} \to \P^1_{C}$ admet une section, sa fibre générique est une quadrique, donc la $C$-variété $X_{C}$ est rationnelle. Pour tout corps $k$ de caractéristique différente de $2$, d'après le théorème de Merkurjev \cite{merkurjev1981norm}, l'application $H^2(k,\mu_{2^n}^{\otimes 2})\to H^2(k,\mu_{2^{n-1}}^{\otimes 2})$ est surjective pour tout $n\geq 1$, et donc l'application $H^3(k,\Z/2) \to H^3(k,\Q/\Z(2))$ est injective. On~en déduit l'injectivité des flèches horizontales dans le diagramme commutatif suivant
\[
\begin{tikzcd}
H^3(R(\P^1),\Z/2) \arrow[r,"\ts \sim"] \arrow[d] & H^3(R(\P^1),\Q/\Z(2)) \arrow[d,->>] \\
H^3_{\on{nr}}(R(X)/R(\P^1),\Z/2) \arrow[r,hook] & H^3_{\on{nr}}(R(X)/R(\P^1),\Q/\Z(2))
\end{tikzcd}
\]
La surjectivité de la flèche du haut suit de la nullité de $H^3(C(\P^1),\Q/\Z(2))$ et d'un argument de restriction-corestriction. La~surjectivité de l'application de droite
est donnée par le théorème \ref{KRS}\eqref{KRSii}. On~en déduit la surjectivité de la flèche de gauche. Combinant ces faits, on voit que toute classe \[
\alpha \in
H^3_{\on{nr}}(R(X)/R,\Z/2) \subset H^3_{\on{nr}}(R(X)/R(\P^1),\Z/2)
\] est l'image d'une classe $\beta \in
H^3(R(\P^1), \Z/2)$.
Soit $\alpha \in H^3_{\on{nr}}(R(X)/R,\Z/2)$. On~fixe un point $m$ de $X(R)$. Sous l'hypothèse que $X(R)$ est semi-algébriquement connexe,
l'image de $X(R)$ dans $\P^1(R)$ est un intervalle de $\P^1(R)$.
Si cet intervalle est tout $\P^1(R)$, le théorème \ref{wittgen} assure que la fibration $X \to \P^1_{R}$ admet une section
et donc que l'espace total est $R$-rationnel. On
supposera que ce n'est pas le cas.
On peut supposer l'intervalle donné par $0 \leq u \leq \infty$,
où $\A^1_{R}=\Spec(R[u]) \subset \P^1_{R}$.

Soit $G\coloneqq \on{Gal}(C/R)$. Comme le corps $C(u)$ est de dimension cohomologique $1$, la~suite exacte longue de cohomologie associée à la suite de $G$-modules \[0\to \Z/2\to (\Z/2)[G]\to \Z/2\to 0\] montre la surjectivité de l'application $H^i(R(u),\Z/2)\to H^{i+1}(R(u),\Z/2)$ donnée par $a\mapsto (-1)\cup a$ pour tout $i\geq 1$. Donc toute classe dans $H^3(R(u),\Z/2)$ s'écrit
comme un cup-produit $(-1,-1,P(u))$ avec $P(u)\in R[u]$. Comme $(-1,r^2+s^2)=0$ dans $\on{Br}(R(u))$ pour tout $r,s\in R(u)$, on peut de plus supposer
\[
P(u) = \prod_{i=1}^n (u+a_{i})
\]
avec tous les $a_{i} \in R$ distincts. La~classe $\alpha$ prend une valeur constante dans $\Z/2= H^3(R, \Z/2)$ sur chaque composante semi-algébrique connexe de $X(R)$, donc sur $X(R)$. Quitte à modifier $\alpha$ par l'image de $(-1,-1,-1)$, on peut supposer que $\alpha$ s'annule sur $X(R)$. Ainsi la fonction $P(u) $ est positive sur $u>0$. Ceci implique $a_{i} \geq 0$ pour tout $i$.

On peut supposer la fibre générique $X_{\eta}$ de $X\to \P^1_R$ donnée par l'annulation d'une forme quadratique $q$ de rang $4$ qui représente $1$.

Soit $f\in R(u)$.
Si la forme quadratique non dégénérée $q$ est une sous-forme de la forme quadratique
$\langle \!\langle -1,-1, f\rangle \!\rangle $, alors cette forme de Pfister admet un zéro sur le corps des fonctions
$R(u)(q)$ de la quadrique définie par $q$, donc est totalement hyperbolique sur
ce corps. Ceci implique que la classe de cohomologie $(-1,-1,f) \in H^3(R(u),\Z/2)$
s'annule dans $H^3(R(u)(q),\Z/2)$.

Pour terminer la démonstration, il~suffit donc de montrer la propriété suivante:
pour tout $a\!\in\! R$ tel que $a \!\geq\!0$, la forme quadratique $q$, de rang $4$,
est une sous-forme de la forme~\hbox{$\!\langle \!\langle -1,-1, u+a \rangle \!\rangle$}, de rang 8.
D'après le théorème \ref{wittgen}, il~suffit de voir
que ceci vaut pour presque toute spécialisation de $u$
en un point $c\in R$.

Soit $a\in R$ tel que $a\geq 0$ et soit $c\in R$.
Si $c+a <0$, alors $c<0$ et donc $c$ n'est pas dans l'image de
$X(R) \to \P^1(R)$. La~forme $q_{c}$ est donc anisotrope.
Comme $q$ représente~$1$ sur $R(u)$, il~en est de même de $q_{c}$ sur $R$,
et $q_{c}$ est isomorphe à $\langle 1,1,1,1\rangle $.
La forme $\langle \!\langle-1,-1, c+a \rangle \!\rangle$ contient la forme
$\langle \!\langle -1,-1 \rangle \!\rangle=\langle 1,1,1,1 \rangle$. Si $c+a>0$, alors la forme $\langle \!\langle -1,-1, c+a\rangle \!\rangle$ sur $R$ est isotrope donc totalement
hyperbolique. Elle représente donc toute forme quadratique sur $R$
de rang au plus égal à $4$.

Voici les modifications à apporter à cette démonstration pour obtenir l'énoncé \eqref{H3nrRpincquadb}. On~considère le diagramme commutatif
\[
\begin{tikzcd}
H^4(R(\P^1),\Z/2) \arrow[r,"\ts \sim"] \arrow[d] & H^4(R(\P^1),\Q/\Z(3)) \arrow[d,->>] \\
H^4_{\on{nr}}(R(X)/R(\P^1),\Z/2) \arrow[r,hook] & H^4_{\on{nr}}(R(X)/R(\P^1),\Q/\Z(3))
\end{tikzcd}
\]
Pour tout corps $k$ de caractéristique différente de $2$, l'application $H^3(k,\mu_{2^n}^{\otimes 3})\to H^3(k,\mu_{2^{n-1}}^{\otimes 3})$ est surjective pour tout $n\geq 1$ par un théorème de Merkurjev, Suslin et Rost,
et donc l'application $H^4(k,\Z/2) \to H^4(k,\Q/\Z(3))$ est injective. Ceci donne l'injectivité des flèches horizontales.
La surjectivité de la flèche verticale de droite est donnée par le théorème \ref{KRS}\eqref{KRSiv}. Le groupe $H^4(R(\P^1),\Q/\Z(3))$ est de 2\nobreakdash-torsion car $C(\P^1)$ est de dimension cohomologique 1. Soit $q$ une forme quadratique de rang 5 sur $R(u)$ représentant 1 et définissant la fibre générique de $X \to \P^1_R$. En utilisant le théorème \ref{wittgen},
et en procédant comme pour la démonstration de \eqref{H3nrRpincquada}, on voit, par spécialisation de $u$ en presque tout point $c\in R$,
que pour tout $a\geq 0$, la forme quadratique
$q$, de rang 5, et donc de rang au plus 8, est une sous-forme de la forme $\langle \!\langle -1,-1,-1, u+a \rangle \!\rangle$, de rang 16.
\end{proof}
\begin{rmk}
La surjectivité de la flèche
$H^{4}(k,\Q/\Z(3)) \to H^{4}_{\on{nr}}(k(q), \Q/\Z(3))$
est aussi établie par Kahn, Rost et Sujatha pour certains types de formes quadratiques de rang $r>5$. Pour celles qui sont de rang au plus 8, la démonstration de la proposition \ref{H3nrRpincquad}(b) s'étend.
\end{rmk}

\begin{rmk}
En combinant \cite[Th.\,1.1]{benoist2026rationality} et \cite[Th.\,4.8]{colliot2024certaines}, on obtient sur $R$ le corps des séries
de Puiseux réelles un
exemple de fibration $X \to \P^1_{R}$ à fibres géométriques des surfaces quadriques intègres,
et un corps $F$ contenant $R$, tels que $X(R)$ soit semi-algébriquement connexe, et que
l'application $H^3(F,\Z/2) \to H^3_{\on{nr}}(F(X)/F,\Z/2)$
ne soit pas surjective.
\end{rmk}

\begin{rmk}
La proposition \ref{H3nrRpincquad}\eqref{H3nrRpincquada} généralise
\cite[Th.\,4.6]{colliot2024certaines}.
Comme observé dans \cite[Rem.\,4.7]{colliot2024certaines},
Benoist et Wittenberg ont donné un critère \cite[Prop.\,5.2]{benoist2020integral1}
assurant que la flèche $H^3(R,\Z/2) \to H^3_{\on{nr}}(R(X)/R,\Z/2)$
est un isomorphisme pour de plus larges classes
de solides $X/R$ pour lesquelles $X(R)$ est semi-algébriquement connexe.
Que ce critère s'applique aux solides fibrés en surfaces quadriques sur
une courbe est un cas particulier de
\cite[Th.\,8.1(i)]{benoist2020integral2}.
\end{rmk}

\begin{cor}\label{H3Rnrdeuxquad}
Soit $R$ un corps réel clos et soit $X \subset \P^5_{R}$ une intersection complète lisse de deux quadriques. Supposons $X(R) \neq \emptyset$. Supposons $X(R)$ semi-algébriquement connexe.
Alors l'application $ H^3(R, \Z/2) \to H^3_{\on{nr}}(R(X)/R, \Z/2)$
est un isomorphisme.
\end{cor}

\begin{proof}
On combine le théorème \ref{theorem3.2} et la proposition \ref{H3nrRpincquad}.
\end{proof}

\begin{prop}\label{fibP1Rrelplus4}
Soit $F$ un corps de caractéristique différente de $2$.
Soit $X\to \P^1_{F}$ une fibration en quadriques de dimension relative $n\geq 4$.
On suppose que $X/F$ est projective et lisse
et que les fibres géométriques singulières ont un seul point singulier
(pinceau de Lefschetz).
On suppose en outre que la fibre générique, qui est définie par
une forme quadratique $q$ sur $F(\P^1)$ de rang $n+2\geq 6 $,
n'est pas semblable à une forme d'Albert anisotrope.
Supposons $X(F)\neq \emptyset$.
Alors l'application $H^3(F,\Z/2) \to H^3_{\on{nr}}(F(X)/F, \Z/2)$ est
un isomorphisme.
\end{prop}

\begin{proof}
Rappelons qu'une forme d'Albert est de rang $6$. Si la forme $q$ est isotrope, alors la fibre générique est une quadrique lisse avec un point rationnel, donc est rationnelle sur le corps $F(\P^1)$,
et $X$ est rationnelle sur $F$. Le résultat est alors clair. Supposons désormais $q$ anisotrope.
Soit $\alpha \in H^3_{\on{nr}}(F(X)/F, \Z/2)$. Comme $q$
n'est pas semblable à une forme d'Albert, d'après
\cite[Th.\,5]{kahn1998unramified},
$\alpha$ est l'image d'une classe $\beta \in H^3(F(\P^1),\Z/2)$.
Les résidus de cette classe en tout point fermé $m \in \P^1_F$
sont dans le noyau de $H^2(F(m), \Z/2) \to H^2(F(X_{m}), \Z/2)$.
Comme chaque $X_{m}$ est défini par une forme quadratique de rang au moins 5,
ce noyau est trivial. Ainsi $\beta$ est dans $H^3(F,\Z/2)$.
\end{proof}

\begin{cor}\label{2quadsurFngeq7}
Soit $F$ un corps de caractéristique différente de $2$, soit \hbox{$n\!\geq\! 7$} et soit $X \!\subset\! \P^n_F$ une intersection complète lisse de deux quadriques. Supposons \hbox{$X(F) \!\neq\! \emptyset$}.
Soit $Y \to \P^1_F$ la fibration en quadriques associée à un $F$-point de $X$ comme dans le théorème \ref{theorem3.2}.
Si $n=7$, supposons que la fibre générique de cette fibration n'est pas définie par une forme d'Albert anisotrope sur $F(\P^1)$. Alors l'application $H^3(F,\Z/2) \to H^3_{\on{nr}}(F(X)/F, \Z/2)$ est un isomorphisme.
\end{cor}
\begin{proof}
On combine le théorème \ref{theorem3.2} et la proposition \ref{fibP1Rrelplus4}.
\end{proof}

\begin{prop}\label{fibP1Rrel3}
Soit $R$ un corps réel clos et soit $X\to \P^1_{R}$ une fibration en quadriques de dimension relative $n= 3$.
On suppose que la variété $X/R$, de dimension~$4$, est projective et lisse
et que chaque fibre géométrique a au plus un point singulier. Supposons que $X(R)$ est semi-algébriquement connexe.
Alors pour tout corps~$F$ contenant~$R$,
l'application $H^3(F,\Z/2) \to H^3_{\on{nr}}(F(X)/F, \Z/2)$ est
un isomorphisme.
\end{prop}
\begin{proof}

Notons $T\subset \P^1(R)$ l'ensemble des points $P$ tels que la fibre $X_P$ soit définie par une forme de rang $4$ sur $R$ anisotrope. Une équation locale de $X\to \P^1_R$ au voisinage d'un point $P\in T$ est de la forme $\sum_{i=0}^3 a_ix_i^2+bx_4^2=0$, où $a_i\in O_{\P^1_R,P}^\times$ et~$b$ est une uniformisante de $O_{\P^1_R,P}$ et $a_i(P)>0$ pour tout $i$.
Pour tout $Q$ proche de~$P$, l'ensemble $X_Q(R)$ est vide ou non selon que $b(Q)>0$ ou $b(Q)<0$. Par contre, soit $P'\in \P^1(R)\setminus T$. On~vérifie sur l'équation locale au voisinage de $P'$ que pour $Q'$ proche de $P'$ la vacuité de $X_{Q'}(R)$ ne dépend pas de $Q'$. On~déduit que la cardinalité~$|T|$ de $T$ est paire. Si on avait $|T|\geq 4$, alors $X(R)$ aurait au moins $2$ composantes semi-algébriquement connexes. Donc $|T|$ est égal à $0$ ou $2$. Si $T$ est vide, alors la projection $X(R) \to \P^1(R)$ est surjective
et le théorème \ref{wittgen} assure que $X \to \P^1_{R}$ a une section
et donc que $X$ est rationnelle sur $R$.
Supposons que $|T|=2$.
On peut donc supposer $T=\{0,\infty\}$.

Soit $\alpha \in H^3_{\on{nr}}(F(X)/F, \Z/2)$. D'après \cite[Th.\,5]{kahn1998unramified} (\cf preuve de la proposition~\ref{H3nrRpincquad}),
$\alpha$ est l'image d'une classe $\beta \in H^3(F(\P^1),\Z/2)$.
Les résidus de cette classe en tout point fermé $m \in \P^1_{F}$
sont dans le noyau de $H^2(F(m), \Z/2) \to H^2(F(X_{m}), \Z/2)$.
Ce noyau est trivial si $X_{m}$ est lisse sur $F(m)$ ou si $X_m$ contient un point
$F(m)$\nobreakdash-rationnel lisse. Soit $m$ un point fermé de $\P^1_{F}$
dont la fibre $X_{m}/F(m)$ n'est pas lisse et est définie par
une forme quadratique de rang $4$ anisotrope.
Alors l'image de~$m$ par la projection $\P^1_{F}\to \P^1_{R}$ est
un point $P\in T$, donc $m=0$ ou $m=\infty$.

Le résidu de $\beta$ en tout point fermé $m$ de $\P^1_{F}\setminus \{0,\infty\}$ est trivial, car ce résidu est
dans le noyau de $H^2(F(m),\Z/2) \to H^2(F(X_{m}),\Z/2)$.

En $0$ et $\infty$, ce résidu est $0$ ou la classe de $(-1,-1)_{F}$.
De plus, par réciprocité \cite[Prop.\,2.2]{rost1996chow}, la somme de ces deux résidus est nulle. Donc
à addition près d'un élément de $H^3(F,\Z/2)$ on a
$\beta=0$ ou $\beta=(-1,-1,u)_{F}$, où $u$ est une fonction rationnelle sur $\P^1_{R}$
de diviseur $(0)-(\infty)$. Pour conclure, il~suffit de montrer que l'image de $(-1,-1,u)_{R}$ dans $H^3(R(X), \Z/2)$ est nulle.

On peut supposer que la fibre générique de $X\to \P^1_R$ est définie par une forme quadratique $q$ sur $R(\P^1)$ de rang $5$ qui représente $1$.
Il suffit de montrer que la forme $q$ est une sous-forme de $\langle \!\langle -1,-1,u\rangle \!\rangle$.
D'après le théorème \ref{wittgen}, ceci vaut si pour presque tout $c\in R$,
la forme $q_{c}$ est une sous-forme de $\langle \!\langle -1,-1,c\rangle \!\rangle$.
Pour presque tout $c>0$ la forme $q_{c}$ de rang $5$ est isotrope et la forme de Pfister $\langle \!\langle-1,-1,c\rangle \!\rangle$ est isotrope donc
totalement hyperbolique. Ainsi $q_{c}$ est une sous-forme de $\langle \!\langle-1,-1,c\rangle \!\rangle$.
Pour presque tout $c<0$, la forme $q_{c}$ de rang $5$ est anisotrope et représente 1, donc
est isomorphe à $ \langle 1,1,1,1,1\rangle $. La~forme $\langle \!\langle-1,-1,c\rangle \!\rangle$ est isomorphe
à $\langle \!\langle-1,-1,-1\rangle \!\rangle$. Ainsi $q_{c}$ est une sous-forme de $\langle \!\langle-1,-1,c\rangle \!\rangle$.
\end{proof}

\begin{prop}\label{fibquaRconnexenqeq3}
Soit $R$ un corps réel clos et soit $X\to \P^1_{R}$ une fibration en quadriques de dimension relative $ n \geq 3$.
On suppose que la variété $X/R$ est projective et lisse
et que chaque fibre géométrique a au plus un point singulier.
Supposons que $X(R)$ est semi-algébriquement connexe.
Pour toute extension $F/R$, l'application $H^3(F,\Z/2) \to H^3_{\on{nr}}(F(X)/F, \Z/2)$ est
un isomorphisme.
\end{prop}
\begin{proof}
Pour $n=3$, c'est la proposition \ref{fibP1Rrel3}. Pour $n\geq 4$,
cela résulte immédiatement de la proposition \ref{fibP1Rrelplus4}, sauf si
$n=4$ et la forme quadratique $q$ définissant la fibre générique
est une forme d'Albert anisotrope. Mais sur le corps $R(\P^1)$,
indice et exposant des éléments du groupe de Brauer coïncident
(tout élément est annulé par passage à $R(\sqrt{-1})$). Ceci implique
qu'un produit tensoriel de deux algèbres de quaternions n'est pas un corps
gauche, et ce dernier fait implique (Albert) que toute forme d'Albert sur $R(\P^1)$ est isotrope.
\end{proof}

\begin{cor}
Soit $X \subset \P^n_R$ une intersection complète lisse de deux quadriques sur un corps réel clos $R$. Supposons $n \geq 6$ et $X(R)$ semi-algébriquement connexe.
Pour tout corps $F$ contenant $R$, l'application
$H^3(F,\Z/2) \to H^3_{\on{nr}}(F(X)/F, \Z/2)$ est un isomorphisme.
\end{cor}
\begin{proof}

On combine le théorème \ref{theorem3.2} et la proposition \ref{fibquaRconnexenqeq3}.
\end{proof}

\begin{rmk}
Dans le cas $n=6$, au moins pour $R=\R$, Hassett, Kollár et Tschinkel \cite{hassett2022rationality}
ont montré bien mieux : toute telle $X \subset \P^6_{R}$ est $R$-rationnelle.
\end{rmk}

Soit $R$ un corps réel clos et $X$ une $R$-variété projective et lisse géométriquement connexe de dimension $n$.
Soit $s\geq 0$ le nombre de ses composantes semi-algébriques connexes.
On a établi dans \cite{colliot1990real} que l'on a
$ H^{i}_{\on{nr}}(R(X)/R,\Z/2) \simeq (\Z/2)^s $
pour $i\geq n+1$. En particulier, si $X(R)$ est semi-algébriquement connexe et $i\geq n+1$, alors l'application $H^i(R,\Z/2)\to H^i_{\on{nr}}(R(X)/R,\Z/2)$ est un isomorphisme.

Rassemblons ici les résultats
obtenus sur l'invariant $ H^{i}_{\on{nr}}(R(X)/R,\Z/2)$
pour les fibrations en quadriques $X$ sur $\P^1_R$
de dimension relative au plus $3$.

\begin{cor}
Soit $R$ un corps réel clos.
Soit $X$ une $R$-variété projective et lisse munie d'un
morphisme $X \to \P^1_{R}$ dont la fibre générique est une quadrique de dimension $d\geq 1$.
Notons $\phi_i : H^{i}(R,\Z/2) \to H^{i}_{\on{nr}}(R(X)/R,\Z/2)$.
\begin{enumeratea}
\item
Pour tout $d$ et $i\leq 1$, $\phi_i$ est un isomorphisme.

\item
Si $d=1$ et $X(R)$ est semi-algébriquement connexe,
alors $X$ est rationnelle et~$\phi_i$ est un isomorphisme pour tout $i$.

\item
Si $d=2$ et $X(R)$ est semi-algébriquement connexe,
$\phi_i$ est un isomorphisme si $i\neq 2$. C'est un isomorphisme pour tout $i$ si toutes les fibres géométriques sont des quadriques avec au plus un point singulier.

\item
Si $d=3$ et $X(R)$ est semi-algébriquement connexe, $\phi_i$ est un isomorphisme si $i\neq 3$. C'est un isomorphisme pour tout $i$ si toutes les fibres géométriques sont des quadriques avec au plus un point singulier.
\end{enumeratea}
\end{cor}

\begin{proof}
D'après le paragraphe précédent, il~suffit de considérer $i\leq d+1$.

\begin{enumeratea}
\item\label{pf:a}
Soit $C\coloneqq R(\sqrt{-1})$. On~a
\[H^1(R,\Z/2)=H^1_{\text{ét}}(X,\Z/2)=H^1_{\on{nr}}(R(X)/R,\Z/2).\]
La première égalité suit du fait que $H^1_{\text{ét}}(X_C,\Z/2)=0$, la $C$-variété $X_C$ étant rationnelle.

\item\label{pf:b}
La rationalité de $X$ est un cas particulier du théorème de Comessatti \cite{comessatti1912fondamenti} (\cf \cite[\S10.2]{colliot2024certaines}). La~bijectivité des $\phi_i$ s'ensuit.

\item\label{pf:c}
Pour $i=3$, voir la proposition \ref{H3nrRpincquad}\eqref{H3nrRpincquada}. La~deuxième partie de l'énoncé suit de la proposition \ref{br-surj}\eqref{br-surjc} et du fait que $X(R)\neq \emptyset$. (Comme on a $d\geq 2$, l'hypothèse additionnelle entraîne que les fibres géométriques sont intègres.)

\item\label{pf:d}
L'application $\phi_2$ est surjective par la proposition \ref{br-surj}\eqref{br-surjc} et elle est injective car $X(R)$ est non vide. Le cas $i=4$ suit de la proposition \ref{H3nrRpincquad}\eqref{H3nrRpincquadb}. La~deuxième partie de l'énoncé suit de l’intégrité des fibres géométriques et de la proposition \ref{fibquaRconnexenqeq3}.\qedhere
\end{enumeratea}
\end{proof}

\subsection{Certaines fibrations en coniques sur \texorpdfstring{$\P^2$}{P2}}
Soit $R$ un corps réel clos. Dans \cite[Prop.\,3.4, Cor.\,3.6]{benoist2020clemens}, Benoist et Wittenberg ont établi la non-rationalité de fibrations en coniques de la forme
\[
y^2+z^2=g(u,v)
\]
lorsque la courbe plane $\Gamma\subset \P^2_R$ d'équation affine $g(u,v)=0$ est lisse de genre au moins $2$. Il est naturel de considérer le cas où $\Gamma$ est singulière et de genre géométrique zéro. Ceci motive la proposition suivante.

\begin{prop}\label{BeWi}
Soit $R$ un corps réel clos et soit $X \to \P^2_{R}$ une fibration en coniques dont un ouvert affine est
défini par une équation
\[
y^2+z^2=g(u,v),
\]
où $g(u,v)$ est un polynôme géométriquement irréductible,
de degré pair tel que $g\geq 0$ sur $\A^2(R)$.
Supposons que la courbe $\Gamma \subset \P^2_R$ d'équation affine $g(u,v)=0$ est à singularités quadratiques ordinaires, est de genre géométrique zéro, et ne coupe la droite
à l'infini $L$ qu'en des points lisses. Pour tout corps $F$ contenant $R$,
l'application
\[
H^3(F,\Q/\Z(2)) \to H^3_{\on{nr}}(F(X)/F,\Q/\Z(2))
\]
est un isomorphisme.
\end{prop}

\begin{proof}
Soit $C\coloneqq R(\sqrt{-1})$. La~variété $X_{C}$ est rationnelle. Si $-1$ est un carré dans $F$, alors $X_F$ est rationnelle et donc pour tout $i$ l'application $H^i(F,\Q/\Z(2)) \to H^i_{\on{nr}}(F(X)/F,\Q/\Z(2))$ est un isomorphisme. Supposons donc que $\sqrt{-1}$ n'est pas un carré dans $F$. Donc $F'\coloneqq F(\sqrt{-1})$ est une extension de corps de $F$ avec groupe de Galois $G=\on{Gal}(F'/F)\cong \Z/2$.

Soient $G_F$ le groupe de Galois absolu de $F$ et $M$ un $G_F$-module galoisien discret de torsion première à $2$. Pour tout entier $i\geq 0$ on a le diagramme commutatif
\[
\begin{tikzcd}
H^i(F,M) \arrow[r]\arrow[d] & H^i(F',M)^G\arrow[d] \\
H^i_{\on{nr}}(F(X)/F,M) \arrow[r] & H^i_{\on{nr}}(F'(X)/F',M)^G
\end{tikzcd}
\]
Les flèches horizontales sont des isomorphismes car $2$ est inversible dans $M$ et la flèche verticale de droite est un isomorphisme car $X_{F'}$ est rationnelle sur $F'$. Il s'ensuit que l'application $H^i(F,M)\to H^i_{\on{nr}}(F(X)/F,M)$ est un isomorphisme.
Il nous reste donc à montrer que l'application \[H^3(F,\Q_2/\Z_2(2))\to H^3_{\on{nr}}(F(X)/F,\Q_2/\Z_2(2))\]
est un isomorphisme.

Soit $\alpha \in H^3_{\on{nr}}(F(X)/F,\Q_2/\Z_2(2))$. D'après Kahn-Rost-Sujatha \cite[Th.\,5]{kahn1998unramified}, il~existe $\beta \in H^3(F(\P^2), \Q_2/\Z_2(2))$ d'image $\alpha \in H^3(F(X),\Q_2/\Z_2(2))$. Comme $X(R)$ est non vide, $X(F)$ est non vide. Soit $P\in X(F)$. En modifiant $\alpha$ par un élément convenable de $H^3(F,\Q_2/\Z_2(2))$, on peut supposer que la spécialisation de $\alpha$ en $P$ est nulle. Comme $X_{F'}$ est rationnelle sur $F'$, l'image de $\alpha$ dans $H^3_{\on{nr}}(F'(X),\Q_2/\Z_2(2))$ est alors nulle. L'extension $F'(X)/F'(\P^2)$ étant transcendante pure, ceci entraîne que l'image de $\beta$ dans $H^3(F'(\P^2), \Q_2/\Z_2(2))$ est aussi nulle.

Soit $\eta$ le point générique de $\Gamma_F$. On~considère le résidu de $\beta$ en un point $\xi$
de codimension $1$ de $\A^2_{F}$ différent de $\eta$. Ce résidu est dans le sous-groupe de $\Br(F(\xi))$ engendré par $(-1,g)$ où on note $g$ la fonction rationnelle
induite par $g$ en $\xi$. Considérons la classe
\[(-1,P(u,v),g(u,v)) \in H^3(F(\P^2), \Z/2) \subset H^3(F(\P^2), \Q_2/\Z_2(2)),\]
où $P(u,v)$ est un polynôme séparable définissant
les $\xi$ autres que $\eta$ pour lesquels le résidu $\partial_{\xi}(\beta)$
est non nul. On~remplace $\beta$ par $\beta - (-1,P(u,v),g(u,v))$.
Cet élément a encore pour image $\alpha$ dans $H^3(F(X),\Q_2/\Z_2(2))$.
On est ramené à supposer que $\beta$ a tous ses résidus
sur $\P^2_{F}$ nuls sauf peut-être $\partial_{\eta}(\beta)$
et $\partial_{L}(\beta)$.

Comme $X_{F'}$ est $F'$-rationnelle, l'image de $\alpha$ dans $H^3_{\on{nr}}(F'(X),\Q_2/\Z_2(2))$ provient de $H^3(F',\Q_2/\Z_2(2))$. Le morphisme $X_{F'}\to \P^2_{F'}$ admet une section rationnelle. Par spécialisation, ceci montre que l'image de $\beta$ dans $H^3(F'(\P^2),\Q_2/\Z_2(2))$ provient de $H^3(F',\Q_2/\Z_2(2))$.

Soit $p \colon S \to \P^2_{R}$ l'éclatement des points singuliers de $\Gamma$. On~considère les complexes de Bloch-Ogus pour $S_F$ et $S_{F'}$:
\[
\begin{tikzcd}
H^3(F(S),\Q_2/\Z_2(2))\arrow[r] \arrow[d] & H^3(F'(S),\Q_2/\Z_2(2)) \arrow[d] \\
{\smash{\bigoplus_{\substack{\gamma\in (S_F)^{(1)}}}} H^2(F(\gamma),\Q_2/\Z_2(1))} \arrow[r] \arrow[d] &
{\smash{\bigoplus_{\substack{\gamma'\in (S_{F'})^{(1)}}}} H^2(F'(\gamma'),\Q_2/\Z_2(1))} \arrow[d] \\
{\smash{\bigoplus_{\substack{m\in (S_F)^{(2)}}}} H^1(F(m),\Q_2/\Z_2)} \arrow[r] &
{\smash{\bigoplus_{\substack{m'\in (S_{F'})^{(2)}}}} H^1(F'(m'),\Q_2/\Z_2)}
\end{tikzcd}
\]
Pour tout $m\in (S_F)^{(2)}$ de corps résiduel $F(m)$ contenant $F'$, on a un isomorphisme $F(m)\otimes_F F'\cong F(m_1)\times F(m_2)$, où $m_1,m_2\in (S_{F'})^{(2)}$ sont les deux $F'$-points d'image $m$. Pour $i=1,2$, l'inclusion naturelle $F(m)\subset F(m_i)$ est une égalité et donc l'application de restriction $H^1(F(m),\Q_2/\Z_2)\to H^1(F(m_i),\Q_2/\Z_2)$ est un isomorphisme pour $i=1,2$. Comme l'image de $\beta$ dans $H^3(F'(S),\Q_2/\Z_2(2))$ provient de $H^3(F',\Q_2/\Z_2(2))$, on conclut que $\partial_m(\partial_\gamma(\beta))=0$ pour tout $\gamma\in (S_F)^{(1)}$ et pour tout $m\in (S_F)^{(2)}$ dans la clôture de $\gamma$ et corps résiduel $F(m)$ contenant $F'$.

Soit $\Delta \subset S$ la désingularisée de $\Gamma$, soit $m\in \Delta_F$ un point fermé et soit $\cl{m}\in \Gamma$ l'image de $m$. Si $\cl{m}$ n'appartient pas à $L_F$ et $\Gamma_F$ est lisse en $\cl{m}$, alors il résulte du complexe de Bloch-Ogus sur $\P^2_F$ que $\partial_m(\partial_\Delta(\beta))=0$. Sinon, comme les points de $\Gamma_F\cap L_F$ et les points singuliers de $\Gamma_F$ sont définis sur $F'$, le corps résiduel $F(m)$ contient $F'$ et donc $\partial_m(\partial_\Delta(\beta))=0$ d'après le paragraphe précédent. On~conclut que $\partial_{\eta}(\beta) \in H^2(F(\Delta),\Q_2/\Z_2(1))$ est non ramifié et donc provient de $H^2(F,\Q_2/\Z_2(1))$. Comme l'image de $\beta$ dans $H^3(F'(\Delta),\Q_2/\Z_2(2))$ est nulle, l'image de $\partial_{\eta}(\beta)$ dans $H^2(F'(\Delta),\Q_2/\Z_2(1))$ l'est aussi. Comme $\Delta_{F}$ est une conique lisse qui a un point rationnel sur $F'$, on conclut que $\partial_{\eta}(\beta)$ peut s'écrire $(-1,\rho)$ avec $\rho \in F^*$.

L'élément $\beta - (-1,\rho, g) \in H^3(F(\P^2),\Q_2/\Z_2(2))$ a encore pour image $\alpha$
dans le groupe $H^3(F(X),\Q_2/\Z_2(2))$ et a tous ses résidus nuls sur $\P^2_{F}$
sauf peut-être au point générique de $L$. D'après \cite[Prop.\,8.6, Rem.\,1.11, Rem.\,2.5]{rost1996chow}, ceci suffit à assurer
que $\beta - (-1,\rho, g)$ appartient à $H^3(F,\Q_2/\Z_2(2))$, et donc
$\alpha$ est dans l'image de $H^3(F,\Q_2/\Z_2(2))$.
\end{proof}

\begin{rmk}
Soit $d$ le degré (pair) de la courbe $\Gamma$. Si $d=2$, la $R$-variété $X$ est une quadrique avec un $R$-point. C'est donc une variété rationnelle. Dans \cite[Th.\,2.2]{benoist2026rationality}, il~est montré que la $R$-variété $X$
est rationnelle si $d=4$, et dans \cite[Th.\,4.5]{benoist2026rationality} qu'elle n'est pas rationnelle si $d\geq 12$.
\end{rmk}

\begin{rmk} En degré cohomologique $2$, on a un résultat général.
Soient $R$ un corps réel clos, $g(u,v) \in R[u,v]$ un polynôme
géométriquement irréductible et $X$ la $R$-variété affine d'équation $y^2+z^2=g(u,v)$. Pour toute extension $F/R$, l'application $\Br(F) \to \Br_{\on{nr}}(F(X)/F)$ est surjective. La~méthode pour établir ce résultat est analogue à celle de la proposition \ref{BeWi} mais plus simple. On~utilise le comportement du groupe de Brauer non ramifié dans une fibration \cite[Ch.\,11]{colliot2021brauer}.
\end{rmk}


% End literal retained input author-body.tex

\begin{thebibliography}{99}
\bibitem{auel2017universal} A. Auel, J.-L. Colliot-Thélène \& R. Parimala - “Universal unramified cohomology of cubic fourfolds containing a plane” , in Brauer groups and obstruction problems , Progress in Math. , vol. 320 , Birkhäuser/Springer, Cham, 2017, p. 29-55
\bibitem{arason1975cohomologische} J. K. Arason - “Cohomologische Invarianten quadratischer Formen” , J. Algebra 36 (1975) no. 3, p. 448-491
\bibitem{bochnak1987geometrie} J. Bochnak, M. Coste \& M.-F. Roy - Géométrie algébrique réelle , Ergeb. Math. Grenzgeb. (3) , vol. 12 , Springer-Verlag, Berlin, 1987
\bibitem{bochnak1998real} J. Bochnak, M. Coste \& M.-F. Roy - Real algebraic geometry , Ergeb. Math. Grenzgeb. (3) , vol. 36 , Springer-Verlag, Berlin, 1998, Translated from the 1987 French original, Revised by the authors
\bibitem{bloch1986algebraic} S. Bloch - “Algebraic cycles and higher $K$ -theory” , Adv. Math. 61 (1986) no. 3, p. 267-304
\bibitem{benoist2026rationality} O. Benoist \& A. Pirutka - “On the rationality of some real threefolds” , J. Inst. Math. Jussieu 25 (2026) no. 3, p. 1555-1577
\bibitem{benoist2020clemens} O. Benoist \& O. Wittenberg - “The Clemens-Griffiths method over non-closed fields” , Algebraic Geom. 7 (2020) no. 6, p. 696-721
\bibitem{benoist2020integral1} O. Benoist \& O. Wittenberg - “On the integral Hodge conjecture for real varieties, I” , Invent. math. 222 (2020) no. 1, p. 1-77
\bibitem{benoist2020integral2} O. Benoist \& O. Wittenberg - “On the integral Hodge conjecture for real varieties, II” , J. Éc. polytech. Math. 7 (2020), p. 373-429
\bibitem{benoist2019intermediate} O. Benoist \& O. Wittenberg - “Intermediate Jacobians and rationality over arbitrary fields” , Ann. Sci. École Norm. Sup. (4) 56 (2023) no. 4, p. 1029-1084, Erratum : Ibid. , 58 (2025), no. 3, p. 829–830
\bibitem{comessatti1912fondamenti} A. Comessatti - “Fondamenti per la geometria sopra le superficie razionali dal punto di vista reale.” , Math. Ann. 73 (1912), p. 1-72
\bibitem{coste1992nash} M. Coste \& M. Shiota - “Nash triviality in families of Nash manifolds” , Invent. math. 108 (1992) no. 2, p. 349-368
\bibitem{colliot1995birational} J.-L. Colliot-Thélène - “Birational invariants, purity and the Gersten conjecture” , in $K$ -theory and algebraic geometry : connections with quadratic forms and division algebras (Santa Barbara, CA, 1992) , Proc. Sympos. Pure Math. , vol. 58, Part 1 , American Mathematical Society, Providence, RI, 1995, p. 1-64
\bibitem{colliot1990real} J.-L. Colliot-Thélène \& R. Parimala - “Real components of algebraic varieties and étale cohomology” , Invent. math. 101 (1990) no. 1, p. 81-99
\bibitem{colliot2016hypersurfaces} J.-L. Colliot-Thélène \& A. Pirutka - “Hypersurfaces quartiques de dimension 3 : non-rationalité stable” , Ann. Sci. École Norm. Sup. (4) 49 (2016) no. 2, p. 371-397
\bibitem{colliot2024certaines} J.-L. Colliot-Thélène \& A. Pirutka - “Certaines fibrations en surfaces quadriques réelles” , 2024, À paraître dans Comment. Math. Helv.
\bibitem{colliot1993groupe} J.-L. Colliot-Thélène \& A. N. Skorobogatov - “Groupe de Chow des zéro-cycles sur les fibrés en quadriques” , -Theory 7 (1993) no. 5, p. 477-500
\bibitem{colliot2021brauer} J.-L. Colliot-Thélène \& A. N. Skorobogatov - The Brauer-Grothendieck group , Ergeb. Math. Grenzgeb. (3) , vol. 71 , Springer, Cham, 2021
\bibitem{colliot1987intersectionsI} J.-L. Colliot-Thélène, J.-J. Sansuc \& P. Swinnerton-Dyer - “Intersections of two quadrics and Châtelet surfaces. I” , J. reine angew. Math. 373 (1987), p. 37-107
\bibitem{cheltsov2024rationality} I. Cheltsov, Y. Tschinkel \& Z. Zhang - “Rationality of singular cubic threefolds over $\mathbb{R}$ ” , Adv. Math. 487 (2026), article ID 110756, 27 pages
\bibitem{delfs1981semialgebraicI} H. Delfs \& M. Knebusch - “Semialgebraic topology over a real closed field. I. Paths and components in the set of rational points of an algebraic variety” , Math. Z. 177 (1981) no. 1, p. 107-129
\bibitem{delfs1981semialgebraicII} H. Delfs \& M. Knebusch - “Semialgebraic topology over a real closed field. II. Basic theory of semialgebraic spaces” , Math. Z. 178 (1981) no. 2, p. 175-213
\bibitem{fulton1975rational} W. Fulton - “Rational equivalence on singular varieties” , Publ. Math. Inst. Hautes Études Sci. (1975) no. 45, p. 147-167
\bibitem{fulton1998intersection} W. Fulton - Intersection theory , Ergeb. Math. Grenzgeb. (3) , vol. 2 , Springer-Verlag, Berlin, 1998
\bibitem{hartshorne} R. Hartshorne - Algebraic geometry , Graduate Texts in Math. , vol. 52 , Springer-Verlag, New York-Heidelberg, 1977
\bibitem{hassett2022rationality} B. Hassett, J. Kollár \& Y. Tschinkel - “Rationality of even-dimensional intersections of two real quadrics” , Comment. Math. Helv. 97 (2022) no. 1, p. 183-207
\bibitem{hassett2021rationality} B. Hassett \& Y. Tschinkel - “Rationality of complete intersections of two quadrics over nonclosed fields” , Enseign. Math. (2) 67 (2021) no. 1-2, p. 1-44, with an appendix by Jean-Louis Colliot-Thélène
\bibitem{kahn2008formes} B. Kahn - Formes quadratiques sur un corps , Cours Spécialisés , vol. 15 , Société Mathématique de France, Paris, 2008
\bibitem{kahn1998unramified} B. Kahn, M. Rost \& R. Sujatha - “Unramified cohomology of quadrics. I” , Amer. J. Math. 120 (1998) no. 4, p. 841-891
\bibitem{kahnsuja1998unramified} B. Kahn \& R. Sujatha - “Unramified cohomology of quadrics. II” , Duke Math. J. 106 (2001) no. 3, p. 449-484
\bibitem{lam2005introduction} T. Y. Lam - Introduction to quadratic forms over fields , Graduate Studies in Math. , vol. 67 , American Mathematical Society, Providence, RI, 2005
\bibitem{merkurjev1981norm} A. Merkurjev - “On the norm residue symbol of degree $2$ ” , Dokl. Akad. Nauk SSSR 261 (1981) no. 3, p. 542-547
\bibitem{merkurjev2008unramified} A. Merkurjev - “Unramified elements in cycle modules” , J. London Math. Soc. (2) 78 (2008) no. 1, p. 51-64
\bibitem{rost1996chow} M. Rost - “Chow groups with coefficients” , Doc. Math. 1 (1996), p. 319-393
\bibitem{scheiderer1994real} C. Scheiderer - Real and étale cohomology , Lect. Notes in Math. , vol. 1588 , Springer-Verlag, Berlin, 1994
\bibitem{scheiderer1996hasse} C. Scheiderer - “Hasse principles and approximation theorems for homogeneous spaces over fields of virtual cohomological dimension one” , Invent. math. 125 (1996) no. 2, p. 307-365
\bibitem{schreieder2021unramified} S. Schreieder - “Unramified cohomology, algebraic cycles and rationality” , in Rationality of varieties , Progress in Math. , vol. 342 , Birkhäuser/Springer, Cham, 2021, p. 345-388
\bibitem{scheiderer2024course} C. Scheiderer - A course in real algebraic geometry—positivity and sums of squares , Graduate Texts in Math. , vol. 303 , Springer, Cham, 2024
\bibitem{sga4II} - Théorie des topos et cohomologie étale des schémas. Tome 2 (M. Artin, A. Grothendieck \& J. L. Verdier, eds.) , Lect. Notes in Math. , vol. 270 , Springer-Verlag, Berlin–New York, 1972, Séminaire de Géométrie Algébrique du Bois-Marie 1963–1964 (SGA 4), avec des contributions de N. Bourbaki, P. Deligne et B. Saint-Donat
\bibitem{witt1937theorie} E. Witt - “Theorie der quadratischen Formen in beliebigen Körpern” , J. reine angew. Math. 176 (1937), p. 31-44
\end{thebibliography}
\end{document}
